\section{Cotangent cohomology and its consequences}\label{sec:cohomology}
\subsection{Direct images and the decisive period derivative}\label{sec:cotangent-images}
The calculations concern every admissible member of the fixed marked
$(3,4)$ period family.  The structure-sheaf direct images and the previously
known Hodge constraints are credited to \cite[Proposition 9.13 and Theorem 9.1]{AC26};
the root-curve version also occurs in \cite[Proposition 6.10]{DGKXv1}.
The source proof of the structure-sheaf calculation is retained below as a
new proof, in accordance with the prescribed T status of this subsection.


\begin{definition}[inv:notation-cotangent; text; Cotangent notation]\label{inv:notation-cotangent}
We use the marked threefold of the question, with its arbitrary admissible
constant $c$.  Coefficients of coherent cohomology are complex.  In this
section $\omega_f=K_X\otimes f^*K_B^{-1}$,
$\mathcal L=\Omega_B^1(\log\Sigma)=\mathcal O_B(1)$,
$\mathcal A=\Omega_X^1(\log D)$, and
$\mathcal E_{\rm cot}=\mathcal A/f^*\mathcal L$.  We write $\gamma,u_{\Lambda}$ also for
their images in $R^1f_*\mathcal O_X$; this convention concerns cohomology
classes and does not identify them with differential-form representatives.
\end{definition}
\uses{The family and marking fixed in Section~\ref{sec:family}.}

\begin{theorem}[inv:A1.0; T; The structure-sheaf input in the precise form used below]\label{inv:A1.0}
The direct images $R^qf_*\mathcal O_X$, for $q=0,1,2$, are
$\mathcal O$, $G_B=\mathcal O\gamma\oplus\mathcal O(-1)$, and $\mathcal O(-1)$.
The inclusion $\mathcal O\gamma\subset G_B$ is a subbundle.  At the cusp,
$(\gamma,u_{\Lambda})$ is a holomorphic frame, and cup product identifies
$\bigwedge^2G_B$ with $R^2f_*\mathcal O_X$.  In the cyclic covering coordinate
$s$, the action on $G_B$ has characters $(1,\zeta^{-a})$, where
$(m,a)=(3,2)$ or $(4,1)$ and $\zeta=e^{-2\pi i/m}$.
These assertions concern full scheme fibres, including $3S_1$ and $4S_2$.
\end{theorem}
\uses{The cyclic and cusp charts; analytic base change \cite[III.4.10, III.4.12]{BS76}; analytic duality \cite[pp.~261, 264--266]{RRV71} and \cite[p.~91, Proposition~2]{RR70}.}

\begin{proof}
\par\smallskip\noindent\textbf{Step 1.}\ We derive the assertion from the local charts and the discriminant;
the same direct images are stated in Proposition 9.13 of
\cite{AC26}, but that proposition is not an input to this proof.
Flatness holds
because a nonzero local parameter on the smooth base is a nonzero divisor
in the local ring of the smooth total space.  The normalization complex
of the cusp has rational surface and curve strata and the cellular
cochain complex of a torus as its constant part.  It gives
$h^q(\mathcal O_W)=(1,2,1)$; the two degree-one classes $\gamma,u_{\Lambda}$ have nonzero
cup product.  To compute the determinant line without assuming a
canonical-bundle formula, we first retain the two finite stabilizers.
Let $\mathfrak B$ be the root curve of orders three and four over $B$,
with root lines $R_1,R_2$, and put $R=R_1R_2^{-1}$.  The presentation
$3R_1=4R_2=H$ gives $\operatorname{Pic}(\mathfrak B)=\mathbb Z R$,
$R_1=R^4$, $R_2=R^3$ and $\mathcal O_{\mathfrak B}(p_0)=R^{12}$;
one may also verify this by removing a line bundle's two stabilizer
characters and descending it to $\mathbb P^1$.

\par\smallskip\noindent\textbf{Step 2.}\ Write $\bar f:X\to\mathfrak B$.  Translation-invariant vertical
fields on the torus charts, extended as logarithmic fields on the cusp
charts, have transition matrices depending only on the base.  They
give a rank-two bundle $V$ on $\mathfrak B$ whose pullback is the
relative logarithmic tangent bundle.  Its fixed field $\xi$ is a
nowhere-zero subline, including in the logarithmic frame at the cusp.
The quotient line $Q=V/\mathcal O\xi$ has automorphy factors $-1/\tau$
and $1/\tau$ for the two finite generators, read directly from the
linear part of the given affine deck transformations.  Thus
$\Delta_{\rm mod}(\tau)^{-1}$ is a meromorphic section of $Q^{12}$.
The discriminant product has no zero in the upper half-plane and
has a simple zero in the cusp coordinate, since
$\tau=(\log t_c)/(2\pi i)+h(t_c)$ with $h$ holomorphic.  Its divisor
therefore proves $Q^{12}=\mathcal O_{\mathfrak B}(-p_0)=R^{-12}$.
The displayed torsion-free Picard group gives
$\det V=Q=R^{-1}$.

\par\smallskip\noindent\textbf{Step 3.}\ The exact logarithmic tangent sequence has quotient
$\bar f^*T_{\mathfrak B}(-p_0)$.  Its determinant, together with
$W=\bar f^*p_0$, gives
$\omega_{X/\mathfrak B}=\bar f^*(\det V)^{-1}$.
The geometric fibres on the root charts have cohomology dimensions
$(1,2,1)$: they are smooth tori at the two finite points, and the
normalization calculation above gives the same result at the cusp.
Analytic base change makes the three root direct images locally free.
Connectedness gives $\bar f_*\mathcal O_X=\mathcal O_{\mathfrak B}$, and fibrewise
Gorenstein duality with its trace pairing now yields
$R^2\bar f_*\mathcal O_X=\det V=R^{-1}$.  The invariant local exponents
of $R^{-1}$ are two at the order-three point and one at the
order-four point.  Removing these characters leaves
$R^{-1}R_1^{-2}R_2^{-1}=R^{-12}$, so its coarse direct image is
$\mathcal O_B(-1)$.  This supplies the degree used below independently of
the asserted direct-image or canonical-bundle formulas in \cite{AC26}.

\par\smallskip\noindent\textbf{Step 4.}\ The same base-change calculation identifies $\gamma,u_{\Lambda}$ with a
frame through the cusp.  At a multiple fibre finite pushforward and the
Reynolds projector identify cohomology with cyclic invariants.  If $b_0$
and $b_1$ have characters $1,\zeta^{-a}$, their invariant frames are
$b_0,s^ab_1$ and $s^a(b_0\wedge b_1)$.  Their products have exactly
these exponents, so the determinant comparison is an isomorphism,
including after reduction modulo $s^m$.  Thus $\deg G_B=-1$.
The invariant flat class $\gamma$ does not enter the holomorphic
two-plane spanned by
$\sigma_1=6\mu\gamma+\tau u_{\Lambda}+w_{\Lambda}$ and
$\sigma_2=\beta\gamma+\mu u_{\Lambda}+\delta$: comparison of the $w_{\Lambda},\delta$
coefficients proves this even on the smooth cyclic covers at their
centres.  It therefore gives the stated subbundle everywhere.
The quotient is $\mathcal O(-1)$ and the extension splits because
$H^1(\mathcal O(1))=0$.  This also specifies which subline in the otherwise
noncanonical splitting is used below.  Coherent Leray now gives
$(h^q(\mathcal O_X))_{q=0}^3=(1,1,0,0)$, since
$H^0(\mathcal O(-1))=H^1(\mathcal O(-1))=H^1(\mathcal O)=0$.
\end{proof}

\begin{lemma}[inv:A1.1; T; Local freeness and the determinant line]\label{inv:A1.1}
Both $\mathcal A$ and $\mathcal E_{\rm cot}$ are locally free, of ranks three and
two.  The vertical vector $\xi=\partial/\partial\zeta_2$ is nowhere
zero as a section of $\mathcal E_{\rm cot}^\vee$.  There is an exact sequence
of holomorphic vector bundles
\begin{equation}\label{inv:eq:A-E-line}
0\longrightarrow N_{\rm det}\longrightarrow\mathcal E_{\rm cot}
 \xrightarrow{\langle\cdot,\xi\rangle}\mathcal O_X\longrightarrow0,
\qquad N_{\rm det}=\det\mathcal E_{\rm cot}=\omega_f(-2S_1-3S_2).
\end{equation}
\end{lemma}
\uses{The local cyclic and normal-crossing charts; inv:A1.0.}

\begin{proof}
\par\smallskip\noindent\textbf{Step 1.}\ At a smooth point of a reduced cusp branch the assertion is the usual
coordinate calculation.  At a double point write $t_c=xy$ with third
coordinate $z$; then
$\mathcal A=R\,d\log x\oplus R\,d\log y\oplus R\,dz$ and its
base subline is $R(d\log x+d\log y)$.  At either triple point write
$t_c=xyz$; then $\mathcal A$ has the three logarithmic coordinate
basis vectors and its base subline is their sum.  In each case that
subline is a direct summand.  These calculations apply on every one
of the three double curves, including their endpoints; no simplicity
of the global divisor is assumed.

\par\smallskip\noindent\textbf{Step 2.}\ At $S_j$ we use the actual free cyclic quotient
$q:J'_j\to N_j$, which is finite and etale on total spaces, and
$t_j=s^{m_j}$ on the base.  On $J'_j$ the logarithmic conormal is
$m_j\,ds/s$, and its coefficient is a unit over $\mathbb C$.  It is again a
direct summand of the logarithmic cotangent bundle.  The comparison
$q^*\mathcal A=\Omega_{J'_j}^1(\log(s=0))$ identifies the quotient
bundles and their exact sequence; faithful flat descent proves the
claim downstairs.  This is not a replacement of $m_jS_j$ by $S_j$.

\par\smallskip\noindent\textbf{Step 3.}\ On the marked torus family $\xi$ is the constant second vertical
vector.  The linear monodromies fix it, and translations preserve it.
On a cusp toric chart it is a nonzero constant vector in the relative
logarithmic tangent lattice of rank two.  A toric coordinate change
$x_i\mapsto c_i(t_c)t_c^{b_i}x_i$ changes $d\log x_i$ by a multiple
of $d\log t_c$, so this constant relative vector descends and is
nowhere zero in that bundle.  It can vanish as an ordinary tangent
vector at a toric stratum; this does not imply vanishing in the
logarithmic bundle.  Hence evaluation is surjective and its kernel
is the determinant line.  Taking determinants of
$0\to f^*\mathcal L\to\mathcal A\to\mathcal E_{\rm cot}\to0$ gives
\[
\det\mathcal E_{\rm cot}=K_X(D)\otimes f^*\mathcal L^{-1}
=\omega_f\bigl(D-f^*\Sigma\bigr)
=\omega_f(-2S_1-3S_2),
\]
which proves \eqref{inv:eq:A-E-line}.
\end{proof}

\begin{proposition}[inv:A1.2; T; The logarithmic direct images by vertical twisting and duality]\label{inv:A1.2}
One has, as holomorphic vector bundles on $B$,
\[
\begin{array}{c|ccc}
q&0&1&2\\\hline
R^qf_*\mathcal E_{\rm cot}&\mathcal O\oplus\mathcal O(-1)&
 \mathcal O^2\oplus\mathcal O(-1)^2&\mathcal O\oplus\mathcal O(-1).
\end{array}
\]
They commute with base change to every full scheme fibre.
\end{proposition}
\uses{inv:A1.0, inv:A1.1; analytic duality \cite[pp.~261, 264--266]{RRV71}.}

\begin{proof}
\par\smallskip\noindent\textbf{Step 1.}\ Put $H=2S_1+3S_2$.  We first compute the line bundle $\mathcal O_X(H)$,
rather than a tensor product of two separately descended lattices.
On the cover at $p_j$, its sections allow poles of order at most $m_j-1$
in $s$.  A fixed character has invariant exponent zero and gains no
pole, whereas the variable character $\zeta^{-a}$ has least allowed
invariant exponent $a-m$.  Thus the inclusion
$R^1f_*\mathcal O\to R^1f_*\mathcal O(H)$ is the identity on $\mathcal O\gamma$ and
changes $s^ab_1$ to $s^{a-m}b_1=t_j^{-1}s^ab_1$ on its quotient.
It makes exactly one upper modification at each finite point.
The same computation applies to $R^2$.  No modification is made at
the cusp or a regular fibre.  The result is
\begin{equation}\label{inv:eq:A-H-images}
f_*\mathcal O(H)=\mathcal O,\qquad
R^1f_*\mathcal O(H)=\mathcal O\oplus\mathcal O(1),\qquad
R^2f_*\mathcal O(H)=\mathcal O(1).
\end{equation}
Here the middle extension splits because $H^1(\mathcal O(-1))=0$.
The local invariant modules are free.  Tensoring them with
$\mathbb C\{t_j\}/(t_j)$ gives invariants of
$s^{-(m_j-1)}\mathbb C\{s\}/(s^m)$, with the indicated surviving character
basis vectors.  The Reynolds idempotent is linear over the base,
so this calculation also proves the comparison with the nonreduced
fibre.  At the cusp the line bundle is trivial.  Cohomology and base
change, applied successively from the top degree, proves the claimed
local freeness and comparison for \eqref{inv:eq:A-H-images}.

\par\smallskip\noindent\textbf{Step 2.}\ For clarity the relative duality map we use is the cup product followed
by the analytic trace, giving a morphism of bundles
\[
R^qf_*(\omega_f\otimes\mathcal O(-H))\longrightarrow
\bigl(R^{2-q}f_*\mathcal O(H)\bigr)^\vee.
\]
Every fibre is a compact Gorenstein surface: it is the Cartier fibre
in a smooth threefold.  Serre duality on that scheme makes the map
an isomorphism on its residue field, also for $m_jS_j$.  The constant
dimensions just computed imply local freeness on both sides, so
Nakayama's lemma makes the bundle map an isomorphism.  Therefore
\begin{equation}\label{inv:eq:A-N-images}
(f_*N_{\rm det},R^1f_*N_{\rm det},R^2f_*N_{\rm det})
=(\mathcal O(-1),\mathcal O\oplus\mathcal O(-1),\mathcal O).
\end{equation}
This is the explicit reduction of analytic relative duality needed
here; the scheme fibre is retained throughout.

\par\smallskip\noindent\textbf{Step 3.}\ The connecting maps of \eqref{inv:eq:A-E-line} are zero on a regular
fibre: its relative cotangent is a trivial bundle and evaluation on
a constant nonzero vector splits.  Source and target direct images
in \eqref{inv:eq:A-N-images} and Statement inv:A1.0 are vector bundles,
so these holomorphic connecting maps, zero on a dense open set, are
zero everywhere.  We obtain extensions in each degree of the listed
$R^qN$ by $R^q\mathcal O$.  Every resulting pair of line-bundle degrees is
zero or minus one; their extension groups are
$H^1(\mathcal O(-1)),H^1(\mathcal O),H^1(\mathcal O(1))$, all zero.  These extensions
give exactly the three asserted bundles.  Applying the already proved
fibre comparisons to the short exact sequences proves base change
for $R^q\mathcal E_{\rm cot}$ as well.  This argument uses the vertical twist
$\mathcal O(H)$ and duality; it does not use the workbench's filtration of
$f_*\mathcal E_{\rm cot}\otimes G_B$.
\end{proof}

\begin{definition}[inv:notation-period-branch; text; The order-four branch]\label{inv:notation-period-branch}
The scalar depends on the fourth root chosen in $t-1=s^4$.  We make
that ambiguity explicit.  We choose the local branch for which, on
the positive $t>1$ arc leading to the distinguished cusp,
$\tau=iy$ with $y>1$, and
\begin{equation}\label{inv:eq:A-s-branch}
\tau=i+\frac{i}{\pi\sqrt{E_4(i)}}s^2+O(s^4),\qquad
\operatorname{Im}\mu_s(0)>0.
\end{equation}
The nonvanishing proved below makes the second condition meaningful;
it fixes the remaining sign of $s$.  Equivalently it is the lower
cusp-regular zero branch of the Weierstrass function in fam:B.1.3.
Replacing $s$ by $i^js$ multiplies the reported derivative by $i^{-j}$.
The workbench instead normalizes $ds/dz(z_2)=1$; its number is obtained
by the chain rule and is not numerically the same unspecified coordinate.
\end{definition}
\uses{The period functions and the first-run clockwise convention.}

\begin{theorem}[inv:A2.1; T; The complete raw period determinant]\label{inv:A2.1}
On the marked upper half-plane,
\begin{equation}\label{inv:eq:A-modular-K}
6\mu_z^2-\tau_z\beta_z
=\tau_z^2\frac{E_4(\tau)^2}{E_6(\tau)}.
\end{equation}
The right side extends holomorphically through the apparent poles.
Its exact orders are two at $z_1$ and zero at $z_2$; it has no other
zero away from their translates.  The logarithmic connecting map
$\kappa_0:f_*\mathcal E_{\rm cot}\to\mathcal L\otimes G_B$ has divisor
$\operatorname{divisor}(\det\kappa_0)=p_1+p_2$.
\end{theorem}
\uses{inv:A1.0; the period transformations \cite[Definition 3.1 and Theorem 3.4(i)]{AC26}; the modular identities stated in the proof.}

\begin{proof}
\par\smallskip\noindent\textbf{Step 1.}\ Consider the holomorphic quadratic differential
$\mathscr Q=6(d\mu)^2-d\tau\,d\beta$.  Direct differentiation of
the two affine transformation laws in AC26 Definition 3.1 gives
\[
g_j^*\mathscr Q=\tau^{-2}\mathscr Q\quad(j=1,2).
\]
In particular $g_2^2$ fixes $\tau$ and sends
\[
\mu\longmapsto1-\tau-\mu,\qquad
\beta\longmapsto\beta+6\tau+12\mu-6,
\]
which leaves $\mathscr Q$ exactly unchanged.  The producing run reports zero symbolic residuals for all three identities in \texttt{a\_calculations.py}.
The kernel of
$\Delta=\mathbb Z/3*\mathbb Z/4\to\mathbb Z/3*\mathbb Z/2=\operatorname{PSL}_2(\mathbb Z)$
is the normal subgroup generated by $g_2^2$.  The squared multiplier
of $\mathscr Q$ is the ordinary modular multiplier
$(c\tau+d)^{-2}$, so it is one on every kernel element.  Two generic
points of the marked cover with the same $\tau$ have the same $t$
and differ by an element of $\Delta$; its modular image fixes a
generic point, hence is the identity.  It follows that
$k=\mathscr Q/(d\tau)^2$ is a single-valued meromorphic function of
$\tau$, transforming as a modular form of weight two.  This kernel
argument is needed: mere equivariance of the separate period functions
would not prove descent.

\par\smallskip\noindent\textbf{Step 2.}\ At $z_1$, write a linearizing parameter $r$ with
$g_1r=\zeta r$, $\zeta=e^{-2\pi i/3}$.  If $\mathscr Q=K_rdr^2$
has first nonzero coefficient of order $n$, its transformation law
at the fixed point reads $\zeta^{n+2}=\rho^{-2}=\zeta$.
Thus $n\equiv2\pmod3$, and $K_r$ vanishes to order at least two.
The map $\tau$ is unramified there by AC26 Theorem 3.4(i).  At
$z_2$, $\tau-i$ has order two and $K_s$ is holomorphic; therefore
$k$ has at most a simple pole at $i$.  Away from these points
$\tau$ is unramified.  At the cusp $\mu$ and $\beta+\tau$ are
holomorphic in $q$, so $k=1+O(q)$.

\par\smallskip\noindent\textbf{Step 3.}\ The modular function $k E_6/E_4^2$ is therefore holomorphic at
$\rho$, at $i$, and at the cusp, as well as everywhere else.
The usual zero statements for $E_4,E_6$ (simple at $\rho,i$,
respectively) are the modular facts used in AC26 Section 3.1.
This function descends to the compact modular curve $X(1)=\mathbb P^1$,
so it is constant.  Its cusp value is one.  This proves
\eqref{inv:eq:A-modular-K} and the exact two local orders, without the
workbench's degree count for the coarse determinant line.

\par\smallskip\noindent\textbf{Step 4.}\ For completeness we translate to the actual coarse lattices.  Choose
cover eigenframes $e_0,e_1$ of $\mathcal E_{\rm cot}$ with characters
$1,\zeta^a$, and $b_0,b_1$ of $G_B$ with characters
$1,\zeta^{-a}$.  Their invariant bases are
$e_0,s^{m-a}e_1$ and $b_0,s^ab_1$.  The logarithmic base frame is
$\ell=dt_j/t_j=m\,ds/s$.  Changing determinants from these frames
to the marked ones gives
\[
\det\kappa_0=\epsilon(t_j)m^{-2}s^{m+2-2a}K_s,
\qquad \epsilon(0)\ne0.
\]
This equality follows by taking the two factors $s/m$ from replacing
$ds$ by $\ell$, the factor $s^{m-a}$ in the source, and the factor
$s^{-a}$ in the target.  Thus at $(m,a)=(3,2)$ the total order in
$s$ is $1+2=3$, and at $(4,1)$ it is $4+0=4$.
Both coarse zeros are simple.  At the cusp, the frame change has
no elementary divisor and $dZ/d\tau$ tends to
$\left(\begin{smallmatrix}0&1\\-1&0\end{smallmatrix}\right)$,
with determinant one.  Formula \eqref{inv:eq:A-modular-K} excludes any
further zero, proving the asserted divisor.
\end{proof}

\begin{theorem}[inv:A2.2; T; The evaluated scalar and the two local leading terms]\label{inv:A2.2}
In the coordinate \eqref{inv:eq:A-s-branch},
\begin{equation}\label{inv:eq:A-mu-value}
\mu_s(0)=i\,C_4,
\qquad
C_4=\frac{\sqrt{2/3}}{\pi E_4(i)^{1/4}}
=\frac{4\sqrt\pi}{3^{3/4}\Gamma(1/4)^2}>0.
\end{equation}
Numerically,
\[
C_4=0.2366093140807852246877051134504496062943051340858243372192243072\ldots.
\]
For $r^3=t$, with $r>0$ on the arc $0<t<1$ and $\tau$ travelling
from $\rho$ to $i$ along the unit circle, set
\[
\lambda=\frac{3e^{5\pi i/6}}{2\pi E_6(\rho)^{1/3}}.
\]
Then $\tau=\rho+\lambda r+O(r^2)$ and
\[
K_r=-\frac{3i\lambda}{2\pi}r^2+O(r^3).
\]
The leading coefficient has real and imaginary parts
\[
\begin{aligned}
\operatorname{Re}K_2&=0.08010221108753827140763650309934451243\ldots,\\
\operatorname{Im}K_2&=0.13874109940222333968402480721280640441\ldots.
\end{aligned}
\]
\end{theorem}
\uses{inv:A2.1; the already identified elliptic section fam:B.1.3; the displayed Ramanujan identities.}

\begin{proof}
\par\smallskip\noindent\textbf{Step 1.}\ Ramanujan's differential relations give
$E_6'(i)=-\pi i E_4(i)^2$ and
$E_4'(\rho)=-(2\pi i/3)E_6(\rho)$.
Using $j=1728E_4^3/(E_4^3-E_6^2)$, we obtain
\[
j(\tau)-1728=-1728\pi^2E_4(i)(\tau-i)^2+O((\tau-i)^3).
\]
Since $j=1728(1+s^4)$, the chosen branch is precisely
\eqref{inv:eq:A-s-branch}.  Formula \eqref{inv:eq:A-modular-K} now gives
\[
K_s(0)=-\frac4{\pi^2\sqrt{E_4(i)}}=6\mu_s(0)^2,
\]
because $\tau_s(0)=0$.  This proves nonvanishing and the first
formula in \eqref{inv:eq:A-mu-value}; the sign is fixed by the stated
coordinate convention.  Similarly
$j(\tau)=-512i\pi^3E_6(\rho)(\tau-\rho)^3+cdots$.
This gives the stated $\lambda$, and substitution into
\eqref{inv:eq:A-modular-K} gives the coefficient $-3i\lambda/(2\pi)$.

\par\smallskip\noindent\textbf{Step 2.}\ There is also a geometric nonvanishing proof independent of this
modular determinant.  Using the section identification proved in
fam:B.1.3, set $x=s^2X,y=s^3Y,t=1+s^4$.  The smooth family becomes
\[
Y^2=X^3-3(1+s^4)^3X+2(1+s^4)^4s^2,
\qquad P=(0,-\sqrt2(1+s^4)^2s)
\]
on the chosen branch.  At $s=0$ its equation is
$Y^2=X^3-3X$ and the point is $(0,0)$.  The first variation of the
curve itself is zero.  In the local coordinate $Y$ at that point,
$dX/(2Y)=dY/(3X^2-3)$, so the Abel integral of $P$ has derivative
$\sqrt2/3$.  The marked period is $-i\Omega_R$, rather than the
positive real period $\Omega_R$: the positive imaginary branch
$\tau=iy$, $y>1$, has elliptic scale $F/s=i/\Omega_R$ at zero.
Its period has zero first derivative.  Division gives
$\mu_s(0)=i\sqrt2/(3\Omega_R)\ne0$.
Using $g_2(\mathbb Z+i\mathbb Z)=(4\pi^4/3)E_4(i)$ and
$g_2(\Omega_R(\mathbb Z+i\mathbb Z))=12$ gives
$\Omega_R=\pi E_4(i)^{1/4}/\sqrt3$, which identifies this value
with \eqref{inv:eq:A-mu-value}.  This also explains why choosing the
positive real period with a positive real $s$ instead would select
the reciprocal branch $\operatorname{Im}\tau<1$.

\par\smallskip\noindent\textbf{Step 3.}\ The producing run reports the following numerical methods and outputs.  The first evaluates
$E_4(i)=1+240\sum_{n=1}^{100}\sigma_3(n)e^{-2\pi n}$ at 100
decimal working digits and inserts it in \eqref{inv:eq:A-mu-value}.
The second integrates
\[
\Omega_R=\frac{\sqrt2}{3^{1/4}}
\int_0^{\pi/2}(1-\tfrac12\sin^2\theta)^{-1/2}\,d\theta,
\qquad C_4=\frac{\sqrt2}{3\Omega_R}.
\]
This is the real period of $Y^2=X^3-3X$ for $dX/(2Y)$; substituting
the three roots $\sqrt3,0,-\sqrt3$ in the elementary elliptic-integral
change of variable proves the formula.  The beta-integral evaluation
of the displayed complete elliptic integral gives the gamma
expression in \eqref{inv:eq:A-mu-value}.  The Fourier and quadrature
answers agree with absolute discrepancy $1.79\cdot10^{-102}$.
As a third check, Newton iteration on the Weierstrass Fourier series
with $j(\tau)=1728(1+s^4)$, using the two roots adjacent to
$(1-\tau)/2$, gives their symmetric derivative.  At $s=10^{-15}$ it
is
\[
i\,0.23660931408078522468770511345048081721971218\ldots,
\]
with discrepancy $3.13\cdot10^{-32}$ and maximum residual below
$2.9\cdot10^{-101}$.  The order-three coefficient was independently
evaluated using the divisor-sum expansion of $E_6$ and its Jacobi
theta-constant formula, with agreement to the working precision.
The source reports that the code, output, and precision accompanied that run.  These are numerical
corroborations, not interval certificates; nonvanishing and the exact
orders were proved analytically above.  For example the elementary
bound $\sigma_3(n)\le n^4$ and $e^{-2\pi}<1/500$ bounds the omitted
Fourier tail after $N$ terms by
$250(N+1)^4 500^{-(N+1)}$; roundoff is separately controlled by
the working precision and the independent calculation.
\end{proof}

\begin{theorem}[inv:A2.4; T; The global rank and Lead 2]\label{inv:A2.4}
The map
$\kappa_1:R^1f_*\mathcal E_{\rm cot}\to\mathcal L\otimes R^2f_*\mathcal O=\mathcal O$
is surjective, and $H^0(\kappa_1):\mathbb C^2\to\mathbb C$ has rank one.
Let $h=a\,d\zeta_1+d\zeta_2$ be the unique global section of
$f_*\mathcal E_{\rm cot}$ with $h(\xi)=1$.  In a genuine invariant target
frame normalized as below,
\[
\kappa_1(h\gamma)(p_2)=-\tfrac14\mu_s(0)=-\tfrac i4 C_4.
\]
The rank-one condition is equivalent to
$a\tau_z+\mu_z\not\equiv0$.  Under the same direct-image and
sheaf-surjectivity hypotheses a rank-zero global map would instead
give $h^{1,2}=2$, $h^{1,1}=3$; that alternative does not occur here.
\end{theorem}
\uses{inv:A1.2, inv:A2.1, inv:A2.2; for the explicitly counterfactual final clause, the residue calculation in inv:A1.3.}

\begin{proof}
\par\smallskip\noindent\textbf{Step 1.}\ Let $d_{ij}(s)$ be the cover Kodaira--Spencer coefficients in the
eigenframes of Statement inv:A2.1 after removing $ds$.
The invariant coarse matrix and its cup-product row have the forms
\[
[\kappa_0]=\begin{pmatrix}A&t_jB\\C&D_0\end{pmatrix},\qquad
[\kappa_1]=(-C,A,-D_0,B),
\]
in the ordered basis
$e_0b_0,s^ae_0b_1,s^{m-a}e_1b_0,e_1b_1$.  Specifically
$A=sd_{00}/m$, $B=s^{1-a}d_{01}/m$,
$C=s^{1-a}d_{10}/m$, $D_0=s^{m+1-2a}d_{11}/m$.
These are holomorphic invariant functions; the row follows by
wedging the connecting class with $b_0,b_1$.
If the row had no unit at a finite point, every entry would vanish
modulo $t_j$ and $\det\kappa_0$ would vanish to order at least two,
contradicting inv:A2.1.  Away from the finite points, including the cusp,
$\kappa_0$ is invertible, which gives a unit in the cup-product row.
Thus $\kappa_1$ is surjective.

\par\smallskip\noindent\textbf{Step 2.}\ The lift $h$ exists and is unique because $f_*\mathcal E_{\rm cot}$ is an
extension of $\mathcal O$ by $\mathcal O(-1)$.  Its invariance under $g_2$ gives
$a(g_2z)=\tau a+\mu$, hence $a(z_2)=1/2$.
In the quotient frame $(\gamma,u_{\Lambda})$ of $G_B$ the connecting map is
\[
\kappa_0(h)=\frac{s}{4}
\bigl((6a\mu_s+\beta_s)\gamma+(a\tau_s+\mu_s)u\bigr)\ell.
\]
Choose the variable eigenframe $b_1$ by Reynolds averaging so that
$\gamma\wedge b_1=v(s)\gamma\wedge u_{\Lambda}$, $v(0)=1$.
Then $s\,v(s)(\gamma\wedge u_{\Lambda})\ell$ is the actual invariant
target frame; the uncorrected expression with $v=1$ need not be
invariant away from zero.  Cup product with $\gamma$ gives at the
central point $-(a\tau_s+\mu_s)/4=-\mu_s/4$, proving rank one.

\par\smallskip\noindent\textbf{Step 3.}\ To check the stated equivalence, there is a second global section
\[
q_E=\sigma_1\otimes u_{\Lambda}+6\sigma_2\otimes\gamma
\quad\text{of }R^1f_*\mathcal E_{\rm cot}.
\]
It is the negative of the relative $(1,1)$ image of the invariant
flat class $\eta=u_{\Lambda}\wedge w_{\Lambda}+6\gamma\wedge\delta$.
It extends at the cusp because the relative logarithmic forms and
$\gamma,u_{\Lambda}$ extend there.  At a finite point it is an invariant
holomorphic tensor on the smooth cover, hence belongs to the four
explicit invariant generators above.  Its cup-product row is zero,
since the two terms are $6\mu_z$ and $-6\mu_z$.
It is independent of $h\gamma$ on every regular fibre (its
$\sigma_1\otimes u_{\Lambda}$ coefficient is nonzero), so these two sections
form a basis of the two-dimensional global section space.
Consequently the global map has rank one precisely when its value
on $h\gamma$, equivalently $a\tau_z+\mu_z$, is not identically
zero.  If its rank were zero, the long exact cohomology sequence of
$0\to\mathcal K\to\mathcal O^2\oplus\mathcal O(-1)^2\to\mathcal O\to0$
would give $h^0(\mathcal K)=2$, $h^1(\mathcal K)=1$.
The residue calculation in inv:A1.3 and Leray would then give exactly
the two counterfactual Hodge numbers stated in the claim.
\end{proof}

\begin{theorem}[inv:A1.3; T; All absolute cotangent direct images]\label{inv:A1.3}
For every admissible $c$, the six direct images are
\begin{equation}\label{inv:eq:A-absolute-table}
\begin{array}{c|cc}
q&R^qf_*\Omega_X^1&R^qf_*\Omega_X^2\\\hline
0&\mathcal O(-2)&\mathcal O(-2)\oplus\mathcal O(-1)\\
1&\mathcal O\oplus\mathcal O(-1)^2&\mathcal O(-2)\oplus\mathcal O(-1)^2\\
2&\mathcal O\oplus\mathcal O(-1)&\mathcal O.
\end{array}
\end{equation}
In particular every listed sheaf is torsion-free.
\end{theorem}
\uses{inv:A1.0--inv:A1.2, inv:A2.1, inv:A2.2, inv:A2.4; the normalized residue sequence; analytic relative duality \cite{RRV71}.}

\begin{proof}
\par\smallskip\noindent\textbf{Step 1.}\ We use Statements inv:A2.1, inv:A2.2 and inv:A2.4 above, whose proofs do not use
\eqref{inv:eq:A-absolute-table}.  They give
$\coker\kappa_0=\mathbb C_{p_1}\oplus\mathbb C_{p_2}$,
surjectivity of $\kappa_1$, and rank one on its global sections.
Write $\mathcal K=\ker\kappa_1$.  The logarithmic conormal sequence
then gives
\[
f_*\mathcal A=\mathcal O(1),\quad
0\to T\to R^1f_*\mathcal A\to\mathcal K\to0,\quad
R^2f_*\mathcal A=\mathcal O\oplus\mathcal O(-1),
\quad T=\mathbb C_{p_1}\oplus\mathbb C_{p_2}.
\]
Let $\nu:\widetilde D=\mathrm{dP}_6\sqcup S_1\sqcup S_2\to D$
be normalization.  The local bases in Statement inv:A1.1 give the exact
residue sequence
\[
0\to\Omega_X^1\to\mathcal A\to\nu_*\mathcal O_{\widetilde D}\to0.
\]
Its three residue direct images are
$\mathbb C_{p_0}\oplus\mathbb C_{p_1}\oplus\mathbb C_{p_2}$,
$\mathbb C_{p_1}\oplus\mathbb C_{p_2}$, and zero: the normalization of $W$ is
rational, and a bielliptic surface has $(h^0,h^1,h^2)(\mathcal O)=(1,1,0)$.
In degree zero the map from $\mathcal O(1)$ is local evaluation with the
nonzero multipliers $1,3,4$.  Its kernel is $\mathcal O(-2)$.

\par\smallskip\noindent\textbf{Step 2.}\ We check the torsion cancellation, rather than discarding it.  At
$p_j$, in the invariant frame $(b_0,s^ab_1)$ of $G_B$, residue of
$(dt_j/t_j)(x b_0+y s^ab_1)$ in $H^1(S_j,\mathcal O)$ is $m_jx[b_0]$.
The residue factors through the base logarithmic line on the covering
chart.  Naturality of the connecting map consequently gives
$\operatorname{res}_1\kappa_0=0$.  The determinant of $\kappa_0$
has a simple zero, so its reduction has rank one, and its image is
exactly the one-dimensional kernel $\mathbb C s^ab_1$ of this residue.
Hence $\coker\kappa_0\to H^1(S_j,\mathcal O)$ is an isomorphism.
Both torsion summands cancel, giving the canonical identification
$R^1f_*\Omega_X^1=\mathcal K$ and the asserted $R^2$.
The rank-one global map
$\mathcal O^2\oplus\mathcal O(-1)^2\to\mathcal O$ has a nonzero constant component
on a trivial summand.  Eliminating all other components by bundle
automorphisms gives
$\mathcal K=\mathcal O\oplus\mathcal O(-1)^2$.

\par\smallskip\noindent\textbf{Step 3.}\ Finally $\Omega_X^2=(\Omega_X^1)^\vee\otimes K_X$.  Relative
duality, with the locally free direct images now established, gives
\[
R^qf_*\Omega_X^2
=\bigl(R^{2-q}f_*\Omega_X^1\bigr)^\vee\otimes K_B.
\]
Substitution of $K_B=\mathcal O(-2)$ gives every entry of the second column.
No Euler-characteristic identity is used to select a Hodge number.
\end{proof}

\subsection{The Hodge table}\label{sec:hodge}

The values of $h^{0,q}$, the vanishing of $h^{p,0}$ for $p>0$, and $h^{1,1}=h^{1,2}+1$ were already obtained in \cite[Theorem 9.1 and Remark 1.3]{AC26}. The evaluated middle value completes that table.

\begin{theorem}[inv:A3.1; T; The complete Hodge table for every admissible parameter]\label{inv:A3.1}
For every $\operatorname{Im}c<-M$, the Hodge numbers are
\[
\begin{array}{c|rrrr}
q\backslash p&0&1&2&3\\\hline
3&0&0&0&1\\
2&0&1&2&1\\
1&1&2&1&0\\
0&1&0&0&0.
\end{array}
\]
\end{theorem}
\uses{inv:A1.0, inv:A1.3; coherent Leray on $\mathbb P^1$ and Serre duality.}

\begin{proof}
Since $B$ is a curve, coherent Leray has no differential and gives
for $F=\Omega_X^p$ the exact sequences
$0\to H^1(B,R^{q-1}f_*F)\to H^q(X,F)\to H^0(B,R^qf_*F)\to0$.
Substitution of \eqref{inv:eq:A-absolute-table}, with
$h^0(\mathcal O)=h^1(\mathcal O(-2))=1$ and all cohomology of $\mathcal O(-1)$ zero,
gives columns $p=1,2$ directly.  Statement inv:A1.0 gives column zero,
and Serre duality gives column three.  Thus the numbers are outputs
of the sheaf computation.  Replacing $\beta$ by $\beta+c$ leaves
every derivative, finite eigencharacter, invariant lattice and
logarithmic residue used above unchanged.  The toric gluing units
change, but their relative logarithmic derivatives still differ only
by base logarithmic forms.  Hence the same computations apply at
each admissible $c$, including every point arbitrarily close to the
strict boundary; no manifold at the excluded boundary is asserted.
\end{proof}

\subsection{The Fr\"olicher sequence}\label{sec:frolicher}

The earlier conditional study of a complex structure on $S^6$ supplies
general Hodge and spectral-sequence constraints in
\cite[Proposition 3.1, Theorem 3.2 and Corollary 3.3]{Ugarte00}.
McHugh gives Bott--Chern and Aeppli formulas in terms of Dolbeault and
Fr\"olicher data \cite[Section 7.1, pp.~24--25]{McHugh18}; Angella records
the spectral-sequence and double-complex descriptions and the resulting
Bott--Chern relations \cite[Figures 1--2 and Section 4]{Angella17}.
The geometric input here is the evaluated Hodge table of inv:A3.1;
the proof of inv:C2.1 uses the structure theorem of \cite{Stelzig21}.

Nondegeneration at $E_1$ and injectivity of the map from $H^{0,1}$ were already established in \cite[Corollary 9.22]{AC26}. The following source calculation determines all the first-differential ranks and degeneration at $E_2$.

\begin{definition}[inv:notation-hodge-input; text; The formal Hodge parameter]\label{inv:notation-hodge-input}
In this section $X$ is the standard marked $(3,4)$ threefold and
$h=h^{1,2}(X)=1$ by inv:A3.1.  We retain $h$ in the formal calculations
to specify exactly which conclusions depend on its evaluated value.
Every assertion involving $h$ is valid for every value
consistent with the Hodge table
\begin{equation}\label{inv:eq:C-hodge-input}
\begin{array}{c|rrrr}
q\backslash p&0&1&2&3\\\hline
3&0&0&0&1\\
2&0&h&h+1&1\\
1&1&h+1&h&0\\
0&1&0&0&0.
\end{array}
\end{equation}
Thus no conclusion below silently uses $h=1$.  The geometric inputs for
this part are the structure-sheaf calculation, the vanishing of
holomorphic forms, and the sphere topology, which give
\eqref{inv:eq:C-hodge-input}; Statements inv:A1.3 and inv:A3.1 determine $h$.
We use $d=\partial+\bar\partial$ on complex-valued
smooth forms and the complex orientation for integration.
\end{definition}
\uses{inv:A1.0, inv:A1.3, inv:A3.1; sphere topology.}

\begin{theorem}[inv:C1.1; T; The complete Fr\"olicher differential in terms of $h$]\label{inv:C1.1}
For a compact complex threefold with \eqref{inv:eq:C-hodge-input} and the
complex de Rham cohomology of $S^6$, the Fr\"olicher spectral sequence
satisfies $E_1\ne E_2=E_\infty$.  Its only nonzero $d_1$ ranks are
\[
\rank d_1^{0,1}=1,\quad \rank d_1^{1,1}=h,\quad
\rank d_1^{1,2}=h,\quad \rank d_1^{2,2}=1.
\]
In particular its two middle rows are exact.
\end{theorem}
\uses{The Hodge table in inv:notation-hodge-input; convergence of the Fr\"olicher spectral sequence to de Rham cohomology.}

\begin{proof}
Here $E_1^{p,q}=H^q(X,\Omega_X^p)$ and
$d_r:E_r^{p,q}\longrightarrow E_r^{p+r,q-r+1}$ are complex-linear.
For $r=2$, any source in the two middle rows has target in one of
$(2,0),(3,0),(3,1)$ or outside the table, except the source $(0,2)$,
which is zero.  Sources in row $q=3$ outside $(3,3)$ are zero, and
$(3,3)$ has no target.  For $r=3$ the only possible targets inside the
table have sources $(0,2)$ or $(0,3)$, both zero.  For $r>3$ there is
no arrow inside the table.  A zero term remains zero on subsequent
pages.  Hence $E_2=E_\infty$.

The filtration on $H^k_{\mathrm{dR}}(X;\mathbb C)$ has no nonzero associated
graded term for $1\le k\le5$.  Every middle term of $E_2$ therefore
vanishes.  Consequently
\[
0\longrightarrow\mathbb C\longrightarrow\mathbb C^{h+1}
\longrightarrow\mathbb C^h\longrightarrow0,
\qquad
0\longrightarrow\mathbb C^h\longrightarrow\mathbb C^{h+1}
\longrightarrow\mathbb C\longrightarrow0
\]
are exact, giving the stated ranks.  The first rank is one, so $E_1$
cannot equal $E_2$.
\end{proof}

\subsection{Bott--Chern and Aeppli cohomology}\label{sec:bc-aeppli}

\begin{theorem}[inv:C2.1; T; Bott--Chern and Aeppli tables as functions of $h$]\label{inv:C2.1}
With
\[
H^{p,q}_{BC}=
\frac{\ker\partial\cap\ker\bar\partial\subset A^{p,q}}
{\im\partial\bar\partial},\qquad
H^{p,q}_{A}=
\frac{\ker\partial\bar\partial\subset A^{p,q}}
{\im\partial+\im\bar\partial},
\]
the complete tables are
\begin{equation}\label{inv:eq:C-BC-A}
\begin{array}{c|rrrr}
q\backslash p&0&1&2&3\\\hline
3&0&0&1&1\\
2&0&h&2h&1\\
1&0&2&h&0\\
0&1&0&0&0
\end{array}_{BC}
\qquad
\begin{array}{c|rrrr}
q\backslash p&0&1&2&3\\\hline
3&0&0&0&1\\
2&0&h&2&0\\
1&1&2h&h&0\\
0&1&1&0&0.
\end{array}_{A}
\end{equation}
In particular, substitution of the value $h=1$ proved in inv:A3.1 gives
the workbench's numerical tables.  The totals are
$(1,0,2,2,2,2,1)$ and $(1,2,2,2,2,0,1)$, respectively.
\end{theorem}
\uses{inv:C1.1; the bounded-double-complex decomposition \cite[Theorem 3 and Proposition 6(1)]{Stelzig21}.}

\begin{proof}
\par\smallskip\noindent\textbf{Step 1.}\ We apply the structure theorem for bounded double complexes over a
field, \cite[Theorem 3]{Stelzig21}, to the actual algebraic double
complex $(A^{\bullet,\bullet}(X),\partial,\bar\partial)$.
The word bounded refers to the bidegrees, not to the dimensions of the
spaces of smooth forms, so that theorem applies here.  It gives a
direct-sum decomposition by isomorphisms of double complexes into
squares and zigzags.  An odd zigzag has one-dimensional total
cohomology.  All odd zigzags in intermediate total degree are therefore
excluded by $b_1=\cdots=b_5=0$.  An odd zigzag of length greater than
one would have at least two vertices of its cohomological total degree.
The extreme diagonals $p+q=0$ and $p+q=6$ contain just one bidegree
each, so the remaining odd zigzags are exactly one dot at $(0,0)$ and
one at $(3,3)$.

\par\smallskip\noindent\textbf{Step 2.}\ An even zigzag of length $2r$ contributes a nonzero $d_r$ to one of
the two Fr\"olicher spectral sequences.  inv:C1.1 and complex conjugation
exclude $r\ge2$ in both sequences.  Thus the remaining summands are
arrows.  The horizontal arrows, with their multiplicities, are
\begin{equation}\label{inv:eq:C-arrows}
\begin{array}{c|c|r}
\text{source}&\text{target}&\text{multiplicity}\\\hline
(0,1)&(1,1)&1\\
(1,1)&(2,1)&h\\
(1,2)&(2,2)&h\\
(2,2)&(3,2)&1.
\end{array}
\end{equation}
The vertical arrows are their reflections in $p=q$, with the same
multiplicities.  The correspondence with spectral-sequence ranks is
\cite[Proposition 6(1)]{Stelzig21}; alternatively one computes
column cohomology on each horizontal arrow directly.

\par\smallskip\noindent\textbf{Step 3.}\ We now evaluate the displayed definitions, rather than use a formula
for the dimensions.  On an arrow, the source is killed by neither
the outgoing differential nor its total differential, whereas its
target is closed under both.  There is no composite
$\partial\bar\partial$.  Hence the target contributes one dimension
to Bott--Chern cohomology, and the source contributes one dimension
to Aeppli cohomology.  On a square, only the top right corner is closed
under both operators, and it is the image of
$\partial\bar\partial$ from the bottom left corner.  The other
three corners of the square are images of $\partial$ or
$\bar\partial$, while the bottom left corner is not in
$\ker\partial\bar\partial$.  Thus a square contributes to neither
cohomology.  A dot contributes once to each.  These computations
commute with direct sums, including the possibly infinite sum of
squares.  Counting the sources and targets in \eqref{inv:eq:C-arrows}
and its reflection gives \eqref{inv:eq:C-BC-A}.

\par\smallskip\noindent\textbf{Step 4.}\ The producing run reports that \texttt{code/c\_double\_complex.py} performs a separate
kernel-and-image computation over $\mathbb Q$ and a source-and-target
count.  It includes an acyclic square with its signs, checks
$\partial\bar\partial+\bar\partial\partial=0$, and verifies the
two answers for $h=0,1,2,3,4,5$.  Its output is
\texttt{outputs/c\_double\_complex.txt}.  This is exact linear
algebra, not a floating-point check and not a determination of $h$.
\end{proof}

\begin{announcement}[inv:C3.3; A; The complete Bott--Chern algebra, for every $h$]\label{inv:C3.3}
The positive-total-degree ideal $N_{BC}\subset H^{\bullet,\bullet}_{BC}(X)$
satisfies
\[
H^{\bullet,\bullet}_{BC}(X)=\mathbb C\cdot1\oplus N_{BC},\qquad N_{BC}^2=0.
\]
Its nonzero bidegrees and dimensions are those of inv:C2.1.  Thus the
square-zero conclusion is independent of the evaluated value of $h$.
\end{announcement}
\uses{inv:C1.1, inv:C2.1; repository inv:C1.2 and inv:C3.1; the graded Leibniz rule.}

The repository argument represents the positive Bott--Chern classes by the arrow targets in the double-complex decomposition. It combines the curvature identity of repository inv:C3.1, the graded Leibniz rule, the first-differential ranks and integration in top degree to make the products vanish.
\par\noindent\emph{proof in the repository} (\texttt{INV}, C3.3).

\subsection{Higher multiplicities and finite covers}\label{sec:multiplicities-covers}

For the original multiplicities, the canonical bundle is credited to \cite[Proposition 9.11]{AC26}; the fundamental group and rational homology of the admissible finite-fundamental-group members are credited to \cite[Theorems 7.17 and 7.22]{AC26}. The statements below retain the source hypotheses for higher multiplicities and for the finite universal covers.

\begin{lemma}[inv:B1.1; T; The relative cohomology comparison]\label{inv:B1.1}
Fix the period family and cusp filling. Let $X'$ be obtained by any
admissible translation twists and multiplicities $m_1=3k$, $m_2=4l$,
where $k,l$ are positive integers. Write $D'=W+S'_1+S'_2$ and
\[
\mathcal E_{\rm cot}'=\Omega^1_{X'}(\log D')/f'^*\Omega^1_B(\log\Sigma).
\]
The identity on the common smooth family induces isomorphisms on $B$
\[
R^qf_*\mathcal O_X\simeq R^qf'_*\mathcal O_{X'},\qquad
R^qf_*\mathcal E_{\rm cot}\simeq R^qf'_*\mathcal E_{\rm cot}'\quad(0\le q\le2).
\]
These isomorphisms identify the logarithmic connecting maps $\kappa_q$.
They also hold when the cusp collar is changed by an admissible translation.
\end{lemma}
\uses{The fixed period and cusp charts; inv:A1.0--inv:A1.2; smooth proper base change \cite[III.4.12]{BS76}.}

\begin{proof}
\par\smallskip\noindent\textbf{Step 1.}\ We construct the comparison at every point. Over $B\setminus\Sigma$
the two families are identified, and a translation acts trivially on
$H^q(F,\mathcal O_F)$ and $H^q(F,\Omega^1_F)$ for a complex torus $F$.
For the second assertion, translation fixes the invariant holomorphic
one-forms, and the first follows either from invariant Dolbeault forms
or the exterior algebra $H^\bullet(F,\mathcal O_F)=\bigwedge^\bullet H^1(F,\mathcal O_F)$.
These are statements about the induced maps on cohomology, including
their holomorphic variation.

\par\smallskip\noindent\textbf{Step 2.}\ At a point with original order $d=3$ or $4$, let
$\widetilde f:T\to\Delta_z$ be the smooth proper torus model with
coarse coordinate $q=z^d$. The new model is its pullback by $z=s^e$,
where $m=de$, followed by the free cyclic quotient of order $m$.
Smooth proper base change identifies the higher direct images upstairs
with the pullbacks of those for $T$. The relative logarithmic cotangent
bundle upstairs is precisely $\Omega^1_{T/\Delta_z}$ pulled back:
the base logarithmic direction is killed in the quotient defining $\mathcal E_{\rm cot}'$.
The derivative of a translation contributes only a multiple of $ds$;
it consequently acts on this quotient as ordinary translation on relative forms.

\par\smallskip\noindent\textbf{Step 3.}\ Average over the cyclic action and choose holomorphic eigenframes.
If a frame for $H^q(\mathcal O)$ or $H^q(\Omega^1)$ has character
$\zeta_d^{-a}$, $0\le a<d$, its original invariant extension is
$z^a\mathcal O_{\Delta_q}$. On the order-$m$ cover its character is
$\zeta_m^{-ea}$, and its invariant extension is
$s^{ea}\mathcal O_{\Delta_q}=z^a\mathcal O_{\Delta_q}$.
Taking invariants is exact in characteristic zero; hence this equality
is an equality of the extensions inside the common punctured-disc
bundle, not only an equality of their ranks or degrees. A translation
twist induces the identity on the cohomology just used and changes none
of these lattices. This proves the comparison at both finite points,
including the full nonreduced fibres.

\par\smallskip\noindent\textbf{Step 4.}\ At the cusp the relative logarithmic forms have the global toric frame
$d\log x_1,d\log x_2$, modulo $d\log q$.
A toric gluing has the form
$x_i\mapsto c_i(q)q^{b_i}x_i$ and therefore changes $d\log x_i$ by
$(q c_i'/c_i+b_i)d\log q$. It induces the identity on the relative
logarithmic quotient. The structure-sheaf comparison is also the identity
on the normalization cochain complex. This verifies the cusp extension
and the assertion about translated collars.

\par\smallskip\noindent\textbf{Step 5.}\ The logarithmic cotangent exact sequences are isomorphic on the smooth
base, so their connecting maps agree there under these comparisons.
Their source and target sheaves are locally free by the local calculations
in Section~\ref{sec:cotangent-images} and the eigenframe calculation above. A holomorphic map between
vector bundles on a reduced curve which is zero on a dense open set is
zero everywhere. Thus the connecting maps agree globally.
\end{proof}

\begin{theorem}[inv:B2.1-canonical; T; Canonical bundles for all coprime multiplicities]\label{inv:B2.1-canonical}
Let $X=X_{3k,4l}(u)$ with admissible translation data as in
\cite[Theorem 1.1 and Proposition 2.3]{VL26},
$\gcd(3k,4l)=1$, and any nonzero $u$ in its construction domain.
Then
\[
K_X=f^*\mathcal O_B(-1)\otimes
\mathcal O_X\bigl((k-1)S_1+(3l-1)S_2\bigr).
\]
\end{theorem}
\uses{inv:B1.1; \cite[Theorem 1.1 and Proposition 2.3]{VL26}; relative duality \cite{RRV71}.}

\begin{proof}
\par\smallskip\noindent\textbf{Step 1.}\ The cited proposition supplies exactly the free cyclic quotients used
in inv:B1.1; its Theorem 1.1 supplies the stated sphere topology for its
admissible choices. For the canonical bundle we check the natural evaluation
$f^*f_*\omega_{X/B}\to\omega_{X/B}$.
The invariant relative determinant has weights $a_1=2k$ and $a_2=l$
under the clockwise generators of orders $3k$ and $4l$. The pullback
of a canonical form is $F(s)\,ds\wedge e$, where the character of
$ds\wedge e$ is $\zeta_m^{1+a}$. Therefore its invariant coefficients
are $s^{m-1-a}$ times an invariant function and a nowhere vanishing
unit. Since the quotient of total spaces is etale, the evaluation
vanishes to order $m-1-a$ on the reduced fibre downstairs. These orders
are $k-1$ and $3l-1$. At the cusp the form
$d\log x_1\wedge d\log x_2\wedge dq$ is a nowhere vanishing canonical
frame in every smooth toric chart, so no cusp component occurs.

\par\smallskip\noindent\textbf{Step 2.}\ By inv:B1.1 and relative duality,
$f_*\omega_{X/B}=(R^2f_*\mathcal O_X)^*=\mathcal O(1)$.
The evaluation divisor just computed therefore gives
$\omega_{X/B}=f^*\mathcal O(1)\otimes
\mathcal O((k-1)S_1+(3l-1)S_2)$. Tensoring by $f^*K_B=f^*\mathcal O(-2)$
proves the formula.
\end{proof}

\begin{announcement}[inv:B1.2; A; All admissible twists have the same Hodge numbers]\label{inv:B1.2}
For every family in inv:B1.1 whose cusp is the one-component hexagonal filling,
the direct images of the absolute cotangent sheaves are
\[
\begin{array}{c|ccc}
 &R^0f'_*&R^1f'_*&R^2f'_*\\\hline
\Omega^1_{X'}&\mathcal O(-2)&\mathcal O\oplus\mathcal O(-1)^2&\mathcal O\oplus\mathcal O(-1)\\
\Omega^2_{X'}&\mathcal O(-2)\oplus\mathcal O(-1)&\mathcal O(-2)\oplus\mathcal O(-1)^2&\mathcal O
\end{array}
\]
and $R^qf'_*\mathcal O_{X'}=(\mathcal O,\mathcal O\oplus\mathcal O(-1),\mathcal O(-1))$.
In particular every admissible $(3,4)$ member with
$|12l_0-4l_1-3l_2|=1$ has the Hodge diamond of Section~\ref{sec:hodge}.
\end{announcement}
\uses{inv:B1.1; the repository residue calculation and relative duality.}

The repository compares the invariant extensions of the logarithmic cohomology bundles under the enlarged cyclic roots. The absolute residue maps then remove the same two finite skyscraper terms, and relative duality supplies the second cotangent row.
\par\noindent\emph{proof in the repository} (\texttt{INV}, B1.2).

\begin{announcement}[inv:B2.1-hodge; A; Hodge numbers for all coprime multiplicities]\label{inv:B2.1-hodge}
Let $X=X_{3k,4l}(u)$ with admissible translation data as in
\cite[Theorem 1.1 and Proposition 2.3]{VL26},
$\gcd(3k,4l)=1$, and any nonzero $u$ in its construction domain.
Then $(h^{0,q})_{q=0}^3=(1,1,0,0)$, $h^{p,0}=0$ for $p>0$,
and $(h^{1,1},h^{1,2})=(2,1)$ for every such pair.
Thus these Hodge numbers do not distinguish the multiplicity families.
\end{announcement}
\uses{inv:B1.2; \cite[Theorem 1.1 and Proposition 2.3]{VL26}.}

The announced direct-image comparison gives the structure-sheaf and cotangent cohomology by Leray on the base. These are the Hodge clauses of the same source statement as the preceding canonical formula.
\par\noindent\emph{proof in the repository} (\texttt{INV}, B2.1).

\begin{lemma}[inv:B3.1; T; Connected cyclic covers and the cusp comparison]\label{inv:B3.1}
For an admissible $(3,4)$ member with
$n=|12l_0-4l_1-3l_2|>1$, let $\pi:Y\to X$ be the universal
cover and $g=f\pi$. It is a finite etale holomorphic map of degree $n$.
Every fibre of $g$ is connected; its finite reductions are bielliptic,
its multiplicities are $3,4$, and its cusp $W_n$ normalizes to a
disjoint union of $n$ copies of $\mathrm{dP}_6$.
The pullback maps induce isomorphisms
\[
R^qf_*\mathcal O_X\simeq R^qg_*\mathcal O_Y,\qquad
R^qf_*\mathcal E_{\rm cot,X}\simeq R^qg_*\mathcal E_{\rm cot,Y}.
\]
\end{lemma}
\uses{The admissible affine and cusp fillings; inv:B1.1; the elementary covering and normalization calculations in the proof; \cite[III.4.12]{BS76}.}

\begin{proof}
\par\smallskip\noindent\textbf{Step 1.}\ We first derive the fundamental-group presentation from the fillings.
The smooth family over the punctured base has zero section, so its
fundamental group is $\Lambda\rtimes\langle x,y\rangle$, where
$x,y$ are the clockwise meridians and act by $A_1,A_2$.
On a finite filling the lifted affine generator has $m_j$th power
translation by $v_j$: its $m_j$ translations add to $v_j$ because
$A_jv_j=v_j$. Thus filling the normal disc imposes
$x^3=t_{v_1}$ and $y^4=t_{v_2}$, with the signs fixed by the
clockwise convention on the first page.

\par\smallskip\noindent\textbf{Step 2.}\ The toric covering space of the cusp tube is simply connected.
Indeed its dense torus has cocharacter lattice $\mathbb Z^3$; a loop may
be perturbed into that torus, and the normal discs of the three rays
$(0,1),(e_1,1),(e_2,1)$ kill a basis of this lattice.
This argument applies also to the infinite fan, since a compact loop
and a homotopy meet only finitely many of its locally finite charts.
The same conclusion holds for the small tube: if a loop in
$|t_c|<\epsilon$ bounds a disc in the toric space, the torus action
$(x_1,x_2,t_c)\mapsto(x_1,x_2,a t_c)$ for a sufficiently small
positive $a$ moves that compact disc into the tube, and the scaling
homotopy of its boundary stays in the tube.
The cusp quotient therefore has fundamental group
$\Lambda/\Lambda_{\rm tor}$, and its toric meridian bounds the
normal disc at fixed $x_1,x_2$. A translated collar makes this
meridian $t_\mu y^{-1}x^{-1}$, where $\gamma(\mu)=l_0$.
Van Kampen consequently kills $\Lambda_{\rm tor}$ and imposes
$xy=t_\mu$.

\par\smallskip\noindent\textbf{Step 3.}\ The normal closure of $\Lambda_{\rm tor}$ in the fibre lattice is
$\ker\gamma$: it contains $\hat w,\hat\delta$ and
$A_1\hat w+\hat w=\hat u$, and all these relations lie in
$\ker\gamma$. The remaining generator $c$, the image of
$\hat\gamma$, is central because $\gamma A_j=\gamma$.
The presentation is thus
$\langle c,x,y\mid c\text{ central},\ xy=c^{l_0},\
x^3=c^{l_1},\ y^4=c^{l_2}\rangle$.
Eliminating $y$ makes it abelian, with relation matrix
\[
 \begin{pmatrix}3&-l_1\\4&l_2-4l_0\end{pmatrix},
 \qquad\det=-p.
\]
Its first Smith divisor is one because its entries include three
and four; hence the group is $\mathbb Z/|p|$.
After imposing $c=1$ the quotient has $x^3=x^4=1$, so it is trivial.
Thus the image of the regular-fibre lattice, containing $c$, generates
the entire group $\mathbb Z/n$. The same is true of each reduced special
fibre. At a finite point the cyclic model is a smooth torus family
over a disc, with a free affine finite-group action. An equivariant
smooth trivialization, obtained by averaging a horizontal connection,
contracts the disc to its centre and descends to a deformation
retraction of the local total space onto $S_j$. Thus the image of
the regular-fibre lattice in the local fundamental group is contained
in the image of $\pi_1(S_j)$; the same holds after mapping to
$\pi_1(X)$. At the cusp $\hat\gamma$ survives in
$\Lambda/\Lambda_{\rm tor}$. The following source normal-crossing calculation contracts the tube onto $W$.

\par\smallskip\noindent\textbf{Step 4.}\ Let $N$ be the small cusp tube. The following normal-crossing contraction is the one supplied in inv:B3.3. In each normal-crossing chart
$t_c=z_1\cdots z_r$, $1\le r\le3$, choose the real vector field
whose $z_i$ component is
$-z_i|z_i|^{-2}/\sum_j|z_j|^{-2}$ off the divisor.
It satisfies $\dot t_c=-t_c$ and has norm at most
$C|t_c|^{1/r}$. A finite partition of unity preserves the first
identity and gives a global field with norm at most
$C'|t_c|^{1/3}$ after shrinking the tube. Extend it by zero on $W$.
Its flow exists for all positive time, and its tail length is
bounded by $3C'\epsilon^{1/3}e^{-T/3}$ after time $T$.
Thus the flow converges uniformly to $W$ and gives a continuous
strong deformation retraction fixing $W$. The same estimate proves
continuity at the divisor and at the endpoint of the retraction.

\par\smallskip\noindent\textbf{Step 5.}\ It identifies their fundamental groups, so the same generator is in the image of $\pi_1(W)$.
The covering-space criterion therefore makes the inverse image of
every fibre connected. Admissibility gives $3\nmid l_1$ and odd $l_2$,
so $\gcd(n,12)=1$. The etale cover preserves fibre multiplicities.
On a finite reduction it cannot kill the nontrivial order-$3$ or
order-$4$ linear holonomy, since the covering degree is prime to it;
the resulting finite torus quotient is again bielliptic.

\par\smallskip\noindent\textbf{Step 6.}\ Normalization commutes with the finite etale base change $Y\to X$.
The pullback of the simply connected normalization $\mathrm{dP}_6$
is its trivial degree-$n$ cover. The double curves and triple points
also pull back etale, and their incidence maps are the lifts of the
original incidence maps. Consequently the dual complex is an actual
degree-$n$ covering of the original real two-torus, rather than merely
a cell complex with the same cell counts. It is connected because
$W_n$ is connected. The normalization resolution for $\mathcal O_{W_n}$
has acyclic components, double curves and point strata. Its global
cochain complex is consequently the cellular cochain complex of this
dual torus, with $n,3n,2n$ cells. Thus its coherent cohomology is
$(1,2,1)$, and pullback from $W$ is an isomorphism over $\mathbb C$:
on a finite cover of real tori it is induced by a full-rank map of
their integral first cohomology lattices.

\par\smallskip\noindent\textbf{Step 7.}\ Modulo $d\log q$, the two toric forms $d\log x_i$ are invariant
under all cusp gluings as computed in inv:B1.1. Hence
$\mathcal E_{\rm cot,X}|_W=\mathcal O_W^2$ and $\mathcal E_{\rm cot,Y}|_{W_n}=\mathcal O_{W_n}^2$,
and the same pullback argument applies to their cohomology.
Over a regular fibre, an isogeny of complex tori induces isomorphisms
on $H^q(\mathcal O)$ and $H^q(\Omega^1)$ over $\mathbb C$.
At a finite fibre, make the cyclic base change. The map is an isogeny
of the smooth torus families and induces isomorphisms of their
cohomology bundles. Taking the same cyclic invariant eigenspaces
descends these isomorphisms, including all nilpotent layers of the
scheme fibre. Grauert base change at the cusp, and this explicit
calculation at the finite points, prove the asserted bundle maps are
isomorphisms on all of $B$.
\end{proof}

\begin{theorem}[inv:B3.2; T; Cotangent direct images and Hodge numbers of the covers]\label{inv:B3.2}
Write $T_n=\mathbb C_{p_0}^{\,n-1}$. Then
\[
\begin{array}{c|ccc}
 &R^0g_*&R^1g_*&R^2g_*\\\hline
\Omega^1_Y&\mathcal O(-2)&\mathcal O\oplus\mathcal O(-1)^2\oplus T_n&\mathcal O\oplus\mathcal O(-1)\\
\Omega^2_Y&\mathcal O(-2)\oplus\mathcal O(-1)&\mathcal O(-2)\oplus\mathcal O(-1)^2&\mathcal O\oplus T_n
\end{array}
\]
The splittings involving $T_n$ are not canonical.
The only nonzero Hodge numbers of $Y$ are
\[
h^{0,0}=h^{0,1}=h^{3,2}=h^{3,3}=1,\quad
h^{1,2}=h^{2,1}=1,\quad h^{1,1}=h^{2,2}=n+1.
\]
The original finite-fundamental-group threefold $X$ itself has the
same Hodge diamond as the simply connected $(3,4)$ member.
\end{theorem}
\uses{inv:B3.1; inv:B1.2 (A, source dependency retained); inv:A1.3; derived relative duality \cite{RRV71}.}

\begin{proof}
\par\smallskip\noindent\textbf{Step 1.}\ inv:B3.1 identifies the logarithmic cotangent sequence on cohomology,
and therefore its connecting maps and direct images, with those of
the original family. The normalized residue sheaf differs only at
the cusp: its degree-zero cohomology there is now $\mathbb C^n$.
The residue of a base logarithmic form is the diagonal vector
$(1,\ldots,1)$, since each cusp component has multiplicity one.
Its cokernel is $T_n$. The degree-one residues at the finite fibres
cancel the same two skyscrapers as in inv:B1.2. The long exact sequence
therefore gives
\[
0\longrightarrow T_n\longrightarrow R^1g_*\Omega^1_Y
\longrightarrow\mathcal O\oplus\mathcal O(-1)^2\longrightarrow0
\]
and the remaining entries of the first row. The extension splits
because $\operatorname{Ext}^1(V,T_n)=H^1(V^*\otimes T_n)=0$
for a vector bundle $V$ on $B$.

\par\smallskip\noindent\textbf{Step 2.}\ For clarity relative duality must now be used in its derived form.
For a sheaf $F$ locally free on the total space, it gives
\[
Rg_*(F^*\otimes\omega_{Y/B})
\simeq R\mathcal Hom_B(Rg_*F,\mathcal O_B)[-2].
\]
For $F=\Omega^1_Y$, the torsion in $R^1g_*F$ contributes
$\mathcal Ext^1_B(T_n,\mathcal O_B)$ in degree two of the right side.
This is a length-$(n-1)$ skyscraper at $p_0$. After tensoring by
$\mathcal O(-2)$ it still has that length. The locally free parts dualize
in complementary degrees, giving exactly the second row; the
extension by the locally free quotient splits for the same reason
as above. Coherent Leray computes the displayed Hodge numbers.
For $X$ with one-component cusp, inv:B1.2 already applies independently
of its fundamental group.
\end{proof}

\begin{theorem}[inv:B3.3; T; Canonical rings and rational topology of the covers]\label{inv:B3.3}
For $X$ and $Y$ as above,
\[
K_Y=\pi^*K_X=\mathcal O_Y(-2S'_2),\quad
\bigoplus_{m\ge0}H^0(Y,-mK_Y)=\mathbb C[U,V],\quad
\deg(U,V)=(1,2).
\]
The same ring holds for $X$.
Over $\mathbb C$, the Betti numbers of $Y$ are
$(1,0,n-1,0,n-1,0,1)$; its Fr\"olicher sequence degenerates at $E_2$.
For every nontrivial character of the deck group, the character
part of Dolbeault cohomology consists of one class in $(1,1)$ and
one in $(2,2)$ and no others.
\end{theorem}
\uses{inv:B3.1, inv:B3.2; the source valuation and monomial argument of inv:B2.2, printed below for $(k,l)=(1,1)$; the marked lattice and affine data.}

\begin{proof}
\par\smallskip\noindent\textbf{Step 1.}\ The canonical formula pulls back through an etale map, and
$4S'_2=g^*p_2$. inv:B3.1 gives $g_*\mathcal O_Y=\mathcal O_B$.
The valuation argument in inv:B2.2 therefore applies unchanged and proves
the ring statement with all multiplication maps.

\par\smallskip\noindent\textbf{Step 2.}\ For the invoked source argument, take $m_1=3$, $m_2=4$, $k=l=1$, $d=6$, and integers $r,s$ with $4r+3s=1$, and put $\mathcal L_{\rm root}=\mathcal O(rS_1+sS_2)$. In the following valuation calculation, $f$ denotes either $f:X\to B$ or $g:Y\to B$, and the $S_j$ denote its reduced finite fibres. The weighted root coordinates are $x$ of weight $3$ and $y$ of weight $4$.

\par\smallskip\noindent\textbf{Step 3.}\ The divisor of $f^*(t/(t-1))$ is $m_1S_1-m_2S_2$.
It follows, retaining this rational function as the comparison, that
$\mathcal L_{\rm root}^{m_2}\simeq\mathcal O(S_1)$,
$\mathcal L_{\rm root}^{m_1}\simeq\mathcal O(S_2)$ and
$\mathcal L_{\rm root}^{m_1m_2}\simeq f^*\mathcal O(1)$.
The canonical formula then gives $-K_X=\mathcal L_{\rm root}^d$.

\par\smallskip\noindent\textbf{Step 4.}\ We first establish the required section algebra for every integer
power of $\mathcal L_{\rm root}$. If a meromorphic function represents a section
of $\mathcal O(aS_1+bS_2)$, it is holomorphic on every regular compact fibre
and is constant there. It descends to the base on the smooth locus.
On a cyclic cover $q=s^m$, the finite pole bound gives a Laurent
expansion bounded below, and invariance leaves only powers divisible
by $m$. Thus descent extends meromorphically across the finite point,
with the exact valuation inequality $m\operatorname{ord}_q(h)+a\ge0$.
The identical argument at the reduced cusp has $m=1$.
Conversely each base function satisfying these bounds pulls back to
an allowed section. Consequently, as subsheaves of meromorphic functions,
\[
f_*\mathcal O(aS_1+bS_2)=
\mathcal O_B\left(\left\lfloor a/m_1\right\rfloor p_1+
\left\lfloor b/m_2\right\rfloor p_2\right).
\]
The cohomology used in this valuation calculation is $h^0(\mathbb P^1,\mathcal O(e))=\max\{e+1,0\}$.

\par\smallskip\noindent\textbf{Step 5.}\ For the algebra, let $y,x$ be the canonical sections of
$\mathcal O(S_1),\mathcal O(S_2)$ under the chosen identifications with powers of
$\mathcal L_{\rm root}$. Their powers $y^{m_1},x^{m_2}$ pull back two independent
sections of $\mathcal O_B(1)$, with zeros at the two distinct points.
For any integer $j$, choose the unique residues
$0\le a_0<m_1$, $0\le b_0<m_2$ for which
$j-m_2a_0-m_1b_0=e m_1m_2$ is integral.
The valuation formula identifies $H^0(\mathcal L_{\rm root}^j)$ with
$y^{a_0}x^{b_0}H^0(\mathcal O_B(e))$. When $e\ge0$ a basis is
\[
y^{a_0+m_1i}x^{b_0+m_2(e-i)},\qquad0\le i\le e;
\]
when $e<0$ there are no sections. These are exactly the monomials of
weighted degree $j$. Their multiplication is ordinary polynomial
multiplication, proving the algebra assertion in every degree.
When $k=l=1$, the weight-$6n$ monomials are generated by $x^2$ and
$y^3$, of regraded degrees $1,2$.

\par\smallskip\noindent\textbf{Step 6.}\ The comparisons of inv:B3.1 are pullback maps and identify the logarithmic
cohomology bundles with the invariant part; the deck group acts
trivially on those bundles. At the cusp its action on the $n$
normalization components is the regular permutation representation.
The quotient by the diagonal has each nontrivial character once.
Thus the first-row additional skyscraper in inv:B3.2 has precisely these
characters, and duality gives the conjugate characters in its
second-row additional skyscraper. All other cohomology is invariant
and equal to that of $X$. Coherent Leray now proves the assertion
about the Dolbeault character spaces.

\par\smallskip\noindent\textbf{Step 7.}\ The de Rham complex and its Fr\"olicher sequence split under the
idempotents of the finite deck group. A nontrivial character has just
the two stated Dolbeault entries. No differential can connect them
or connect either to another nonzero entry, so both survive.
The invariant summand is the de Rham complex of $X$ by finite
etale descent and averaging. We compute its rational topology here
from the actual attachment maps; the integral homology assertion
of AC26 is not used as an input. In the remainder of this proof
all homology and cohomology coefficients are $\mathbb Q$.

\par\smallskip\noindent\textbf{Step 8.}\ Let $N$ be a small closed cusp tube, $M=\partial N$, and
$U=X\setminus\operatorname{int}N$. The space $U$ is covered by
two pieces retracting onto $S_1,S_2$, whose intersection retracts
onto a regular fibre $F$. The retractions are obtained from the
equivariant smooth torus trivializations used in inv:B3.1. The maps
$F\to S_j$ are the central finite covering maps; translation of
their origins induces the identity on their rational cohomology.
Finite-cover transfer identifies their pullbacks with the invariant
subspaces for $A_j$. Put
\[
 \psi_1=2u_{\Lambda}+w_{\Lambda}+3\delta,\qquad \psi_2=u_{\Lambda}+w_{\Lambda}+2\delta,
 \qquad \eta=u_{\Lambda}\wedge w_{\Lambda}+6\gamma\wedge\delta.
\]
Direct substitution in $A_1,A_2$ shows that the invariant
degree-one spaces are $\langle\gamma,\psi_j\rangle$ and the
invariant degree-two spaces are
$\langle\gamma\wedge\psi_j,\eta\rangle$.
Their dimensions are two: the nontrivial eigenvalues of $A_j$ on
the complementary plane are the inverse primitive third or fourth
roots, so the exterior-square fixed space has exactly the two
displayed generators.
In the transfer-normalized dual homology bases, the two
Mayer--Vietoris maps are therefore the following actual matrices:
\[
 \alpha_1=\begin{pmatrix}
 1&0&0&0\\0&2&1&3\\-1&0&0&0\\0&-1&-1&-2
 \end{pmatrix},\qquad
 \alpha_2=\begin{pmatrix}
 2&1&3&0&0&0\\0&0&6&1&0&0\\
 -1&-1&-2&0&0&0\\0&0&-6&-1&0&0
 \end{pmatrix}.
\]
The second matrix uses the ordered wedge basis
$(\hat\gamma\hat u,\hat\gamma\hat w,\hat\gamma\hat\delta,
\hat u\hat w,\hat u\hat\delta,\hat w\hat\delta)$.
Both matrices have rank three, and
$\ker\alpha_1=\mathbb Q\nu_0$, with
$\nu_0=-\hat u-\hat w+\hat\delta$.
Consequently
\[
 H_1(U)=\mathbb Q,\qquad
 0\longrightarrow\operatorname{coker}\alpha_2=\mathbb Q
 \longrightarrow H_2(U)
 \xrightarrow{\partial}\mathbb Q\nu_0\longrightarrow0.       \tag{B.1}
\]
The generator $\Theta\in H^1(U)$ is normalized by
$\Theta|_F=\gamma$. The finite affine generators have
$3\bar g_1=l_1c$, $4\bar g_2=l_2c$ in rational homology, where
$c$ is the image of $\hat\gamma$ and $\Theta(c)=1$.
For the toric meridian $a_0=t_\mu y^{-1}x^{-1}$ this gives the
explicit collar value
\[
 \Theta(j_*a_0)=l_0-l_1/3-l_2/4=p/12.                 \tag{B.2}
\]

\par\smallskip\noindent\textbf{Step 9.}\ The contraction of the cusp tube used here was established in the proof of inv:B3.1 by the source normal-crossing vector-field calculation.

\par\smallskip\noindent\textbf{Step 10.}\ The constant-sheaf normalization resolution of $W$ has strata
$\mathrm{dP}_6$, three copies of $\mathbb P^1$, and two points.
Its row of degree-zero cohomology is
\[
 \mathbb Q\xrightarrow{0}\mathbb Q^3
 \xrightarrow{(a,b,c)\mapsto(a-b+c,a-b+c)}\mathbb Q^2.
\]
The strata have no degree-one cohomology. This resolution therefore
gives $H_1(N)=H_1(W)=\mathbb Q^2$, $H_4(N)=\mathbb Q[W]$ and $H_5(N)=0$;
it also gives $e(W)=6-3\cdot2+2=2$.
Poincare--Lefschetz duality for the compact smooth manifold with
boundary $(N,M)$ now shows that $H_1(M)\to H_1(N)$ and
$H_2(M)\to H_2(N)$ are surjective. For degree one the next
relative group is $H_1(N,M)=H^5(N)=0$.
For degree two the map to $H_2(N,M)=H^4(N)$ is the intersection
pairing with $[W]$, which is zero because $W$ is the divisor of
the global function $t_c$ on $N$. The Wang sequence for the
boundary mapping torus, with $\rank(M_0-I)=2$, gives
$\dim H_1(M)=3$. Its kernel to $H_1(N)$ is exactly
$\mathbb Q a_0$: this meridian is nonzero under projection to the base
circle and bounds its normal disc in $N$.

\par\smallskip\noindent\textbf{Step 11.}\ We exhibit two classes in
$K_2=\ker(H_2(M)\to H_2(N))$.
For $\nu=\hat w,\hat\delta$, let $\xi_\nu$ be the two-torus
obtained by translating the toric meridian around the angular
circle $S^1_\nu$. In the chart $t_c=z_0z_1z_2$ corresponding to
the rays $(0,1),(e_1,1),(e_2,1)$, these are the boundaries of
$D^2_{z_0}\times S^1_{z_i}$ with the other coordinate fixed
and the radii chosen so that $|t_c|=\epsilon$ on the boundary.
They therefore bound in $N$ with the literal toric meridian;
no choice of an unrecorded collar lift is being made.
Cutting the meridian into the two arcs of the cover defining
(B.1), the two cylinders have end circles differing by
$(I-A_1^{-1})\nu$. With this orientation of the Mayer--Vietoris
boundary, matrix multiplication gives
\[
 \partial j_*\xi_{\hat w}=-\nu_0,\qquad
 \partial j_*\xi_{\hat\delta}=0.                       \tag{B.3}
\]
The affine translation of either endpoint has no effect on its
homology class, which explains why this boundary formula is the
same for every translation twist.

\par\smallskip\noindent\textbf{Step 12.}\ Choose $\mathcal A\in H^2(U)$ restricting to $\eta$ on $F$;
it exists by Mayer--Vietoris because the two transfer-normalized
classes $\eta$ have equal restrictions. Fibrewise translation by
$S^1_{\hat\delta}$ defines a global smooth action on $U$:
$A_j\hat\delta=\hat\delta$, and all affine translations commute
with this circle. Slant product of its action map with the circle
defines $\iota:H^2(U)\to H^1(U)$. Naturality on $F$ gives
$\iota\mathcal A|_F=\iota_{\hat\delta}\eta=-6\gamma$.
The restriction $H^1(U)\to H^1(F)$ is injective by (B.1), so
$\iota\mathcal A=-6\Theta$. By the definition of slant product,
\[
 \langle\mathcal A,j_*\xi_{\hat\delta}\rangle
 =\langle\iota\mathcal A,j_*a_0\rangle=-p/2.           \tag{B.4}
\]
We may add the class pulled back from $(\mathbb Q\nu_0)^*$ under
the boundary map in (B.1) so that
$\langle\mathcal A,j_*\xi_{\hat w}\rangle=0$;
this does not change its restriction to $F$ or (B.4).
In the coordinates on $H_2(U)$ given by $\partial/\nu_0$ and
this normalized pairing, the two actual attachment cycles have
matrix
\[
 \bigl[j_*\xi_{\hat w}\ \ j_*\xi_{\hat\delta}\bigr]
 =\begin{pmatrix}-1&0\\0&-p/2\end{pmatrix}.             \tag{B.5}
\]
It is invertible over $\mathbb Q$ because admissibility implies
$p\ne0$. Thus the image of $K_2$ already fills $H_2(U)$.
Likewise (B.2) says the image of $\ker(H_1(M)\to H_1(N))$
is nonzero in $H_1(U)$. Using the surjectivity to $N$ just proved,
Mayer--Vietoris for $X=U\cup N$ now gives $H_1(X)=H_2(X)=0$.
Every other fibre has Euler number zero, so additivity over the
three special fibres and the smooth torus bundle gives $e(X)=2$.
Poincare duality then gives $b_4=b_5=0$ and $b_3=0$.
This proves the asserted rational homology of $X$ for all admissible
twists; the only changing determinant in (B.2)--(B.5) is a nonzero
constant multiple of $p$. The producing run reports that the integer matrices and all exterior-power identities were checked exactly in \texttt{code/b\_topology.py},
with output \texttt{outputs/b\_topology.txt}.

\par\smallskip\noindent\textbf{Step 13.}\ The argument of inv:C1.1 using this rational topology and the same Hodge
table now gives degeneration of the invariant summand at $E_2$.
Summing it with the nontrivial character spaces proves the Betti
numbers and the degeneration statement for $Y$. In particular
$e(Y)=2n$, as also follows from the covering degree.
\end{proof}

\subsection{Metric consequences}\label{sec:metric-announcements}

\begin{announcement}[inv:C3.5; A; Metric obstructions that do not require $h$]\label{inv:C3.5}
The threefold $X$ admits no balanced, strongly Gauduchon, K\"ahler,
locally conformally K\"ahler, or Hermitian-symplectic metric.
\end{announcement}
\uses{The fibration and sphere topology; Stokes’ theorem.}

The source uses the nonzero semipositive exact form $f^*\omega_{FS}$ and the positive four-form supplied by a balanced or strongly Gauduchon metric. The remaining exclusions use $b_1=b_2=0$ and the same integral obstruction.
\par\noindent\emph{proof in the repository} (\texttt{INV}, C3.5).

\begin{announcement}[inv:C3.7; A; The fibre obstruction to SKT and astheno-K\"ahler metrics]\label{inv:C3.7}
The threefold $X$ admits no pluriclosed (SKT) metric and no
astheno-K\"ahler metric.  The geometric input is the value $h=1$
and the global Dolbeault class $q$ in inv:A3.2, whose restriction to a
regular fibre is $\eta=u_{\Lambda}\wedge w_{\Lambda}+6\gamma\wedge\delta$.
More generally, it is enough that the real image of
$H_A^{1,1}(X;\mathbb R)\to H_A^{1,1}(F;\mathbb R)$ be contained in
$\mathbb R\eta$ for one regular fibre $F$.
\end{announcement}
\uses{inv:A3.1; the global Dolbeault lift in repository inv:A3.2; inv:C2.1.}

The repository restricts the prospective metric class to a regular torus fibre. Its real Aeppli image lies in $\mathbb R\eta$, while the complex-oriented square of $\eta$ is $-12$; translation averaging would give a positive class with positive square.
\par\noindent\emph{proof in the repository} (\texttt{INV}, C3.7).

\subsection*{Correction concerning relative differentials}\label{sec:relative-correction}
The locators are \cite[Remark 1.3, printed p.~5; Remark 9.21, printed p.~81]{AC26}.
\emph{Paraphrase of the printed assertions:} the former attributes torsion
along the singular fibres to the relative differential sheaves without
separating their degrees; the latter includes the reduced cusp $W$ in the
torsion support of $\Omega^1_{X/B}$. This is a paraphrase, not a quotation:
the requirement not to reproduce the prose of \cite{AC26} prevents the
separately requested literal printed sentence from being supplied.
The correction is local and does not alter the preceding credits.


\begin{lemma}[inv:A4.1; T; All local modules needed to distinguish degrees one and two]\label{inv:A4.1}
Let $R=\mathbb C\{x,y,z\}$.  For $t=xy$,
\[
\Omega_{R/\mathbb C\{t\}}^1=R^3/R(y,x,0)
\simeq(x,y)\oplus R,
\]
and
\[
\Omega_{R/\mathbb C\{t\}}^2\simeq
(x,y)\oplus R/(x,y),\qquad
\Omega_{R/\mathbb C\{t\}}^3\simeq R/(x,y).
\]
For $t=xyz$,
$\Omega^1=R^3/R(yz,xz,xy)$ is torsion-free of rank two and is
not free at the triple point.  If $I=(xy,xz,yz)$, there is an exact
sequence
\[
0\to T\to\Omega^2\xrightarrow{\wedge dt} I\to0,
\qquad
T=\frac{R u\oplus R v}{(xu,\ y(u-v),\ zv)},
\quad IT=0,
\]
and $\Omega^3=R/I$.  Finally, for $t=s^m$ in
$R=\mathbb C\{s,y_1,y_2\}$,
\[
\begin{aligned}
\Omega^1&=R^2\oplus R/(s^{m-1})\,ds,\\
\Omega^2&=R\,dy_1\wedge dy_2
 \oplus\bigl(R/(s^{m-1})\bigr)^2,\\
\Omega^3&=R/(s^{m-1})\,ds\wedge dy_1\wedge dy_2.
\end{aligned}
\]
\end{lemma}
\uses{The presentations of ordinary relative Kähler differentials for $t=xy$, $t=xyz$, and $t=s^m$.}

\begin{proof}
\par\smallskip\noindent\textbf{Step 1.}\ The degree-one presentation is the quotient by $dt$; higher degrees
are its exterior powers.  At a double point the transverse quotient
$R^2/R(y,x)$ maps isomorphically to $(x,y)$ by
$(a,b)\mapsto xa-yb$ after an immaterial change of sign of one
basis vector.  More explicitly one may use
$(a,b)\mapsto xa-yb$ with relation $(y,x)$; its kernel is precisely
that relation, since $x,y$ are relatively prime.  In degree two,
with basis $A=dx\wedge dy,B=dx\wedge dz,C=dy\wedge dz$, the
relations are $xA=yA=0$, $yB+xC=0$.  This gives the displayed
torsion summand and the ideal $(x,y)$; wedging once more gives
$R/(x,y)$.

\par\smallskip\noindent\textbf{Step 2.}\ At the triple point the map from ordinary to logarithmic relative
one-forms sends $(a,b,c)$ to $(xa-zc,yb-zc)$.
Its kernel is precisely $R(yz,xz,xy)$: the two equations imply
$c=xyh$, $a=yzh$, $b=xzh$ by relative primality of the coordinates.
Thus it injects into $R^2$, proving torsion-freeness.  Its fibre at
the maximal ideal has dimension three although its rank is two,
so it is not free.  The same minimal-generator argument works at
the double locus.

\par\smallskip\noindent\textbf{Step 3.}\ In degree two, wedging with $dt$ is the row
$(xy,-xz,yz)$ in the basis $(A,B,C)$.
Its syzygies are generated freely by
$u=(z,y,0)$ and $v=(0,y,x)$.
To check generation, an equation $xyA-xzB+yzC=0$ implies
$x\mid C$, say $C=xC_1$; then $yA=z(B-yC_1)$ implies
$A=zA_1$, $B=y(A_1+C_1)$.  Thus the syzygy is $A_1u+C_1v$.
The three relations from $dt\wedge(dx,dy,dz)$ are
$-xu,y(u-v),zv$, proving the stated presentation of the kernel.
The ideal $I$ kills it by those relations.  Conversely the quotient
$I$ is torsion-free, so this kernel is exactly the torsion of
$\Omega^2$.  The presentation immediately gives $\Omega^3=R/I$.
All matrix identities have also been checked over $\mathbb Z[x,y,z]$ by
the supplied exact program.  For $t=s^m$ the sole degree-one
relation is $ms^{m-1}ds=0$; $m$ is a unit and taking exterior
powers gives the three formulas.
\end{proof}

\begin{remark}[inv:A4.2; T; The claimed degree-one torsion along the cusp]\label{inv:A4.2}
The assertion in AC26 Remark 9.21, PDF page 81, that
$\Omega_{X/B}^1$ has torsion along $W\cup S_1\cup S_2$ is false.
Its torsion is supported precisely on $S_1\cup S_2$; its degree-one
cusp modules are torsion-free but not locally free on the double
locus.  The plural statement in Remark 1.3, PDF page 5, must be
qualified by degree: $\Omega^2_{X/B}$ does have torsion on the cusp
double locus.  The workbench's degree-one correction is correct;
it must not be extended to degree two.
\end{remark}
\uses{inv:A4.1; \cite[Remarks 1.3 and 9.21]{AC26}.}

\begin{proof}
The producing run reports that both disputed sentences were checked on rendered pages of the original PDF. The present source check uses the supplied transcription at the exact locators above.  The
double and triple charts in inv:A4.1 exhaust the singular strata of
$W$.  They prove absence of degree-one torsion at the cusp and
exhibit nonzero degree-two torsion there.  At a multiple fibre the
finite etale comparison
$q^*\Omega^1_{N_j/B}=\Omega^1_{J'_j/B}$ is an isomorphism,
and its last summand in inv:A4.1 is nonzero for $m=3,4$.
Faithful flatness descends both its existence and its support.
There are no other critical fibres and the ordinary relative
cotangent is locally free on their complement.  Thus the stated
support is exact.  The erroneous sentence is at transcription
line 4330 (the final paragraph of Remark 9.21); the workbench's
correct degree-one replacement occupies lines 980--1006 of
\texttt{hodge\_frolicher.tex}.  No numerical Hodge value in that
workbench is refuted by this calculation.
\end{proof}

\section{Tangent cohomology and local deformations}\label{sec:tangent}
\subsection{The standard family: the known Kuranishi line}\label{sec:tangent-known}
The following statement is credited to \cite[Theorem 1.3 and Theorem 6.12]{DGKXv1}.
The first source theorem gives the Kuranishi and miniversality conclusions;
the tangent-cohomology theorem gives the displayed dimensions.
The supplied re-derivation is in Appendix~\ref{app:tangent-rederivation}.


\begin{cited}[fam:A.5.2; K; Tangent cohomology and the Kuranishi line]\label{fam:A.5.2}
For every admissible $c$, the tangent cohomology of the actual compact
member is
\[
 \bigl(h^0(X_c,T_{X_c}),h^1(X_c,T_{X_c}),h^2(X_c,T_{X_c}),
 h^3(X_c,T_{X_c})\bigr)=(1,1,1,0).
\]
Its Kuranishi germ is a smooth reduced disc. The family of fam:A.5.1 is
miniversal at every parameter: its Kodaira--Spencer map is an isomorphism,
and every sufficiently small deformation is induced from this family.
\end{cited}
\uses{The compact family fam:A.5.1 and its nonzero Kodaira--Spencer class fam:A.4.2; \cite[Theorems 1.3 and 6.12]{DGKXv1}.}

\subsection{Every integral Picard twist}\label{sec:tangent-twists}

\begin{definition}[cdp:notation-integral-twists; text; The integral Picard normalization]\label{cdp:notation-integral-twists}
In this subsection $X=X_c$ is the period construction of \cite{AC26}, compared with the quotient construction of \cite{E26}, with admissible parameter $c$, standard clockwise affine vectors
$v_1=(1,2,-4,0)^t$, $v_2=(-1,-3,3,0)^t$, and untwisted cusp collar. Thus $p=12l_0-4l_1-3l_2=-1$. We write $\mathfrak B=\mathbb P(3,4)$, $\varpi:\mathfrak B\to B$ for its coarse map, and $\bar f:X\to\mathfrak B$ for the root factorization. Let $\mathcal T=T_X(-\log W)$ and let $\nu:\mathrm dP_6\to X$ normalize $W$. The root generator is fixed by
\[
 M_n=\bar f^*\mathcal O_{\mathfrak B}(n)
     =\mathcal O_X(n(S_1-S_2)),\quad n\in\mathbb Z.
\]
The exponential isomorphism $H^1(X,\mathcal O_X)\to\Pic(X)$ identifies $\Pic(X)$ with $\mathbb C$. We scale its coordinate so that $M_z=\exp(z\gamma^{0,1})$ and $M_1=\mathcal O_X(S_1-S_2)$. The sign of this equality is fixed by the clockwise transgression $12\gamma\mapsto-p=1$. This is checked by evaluating the two finite lifts on $v_1,v_2$: their integral periods are $4l_1,3l_2$, while the cusp contributes $-12l_0$, whose sum is $-p$. The exponential sequence therefore gives $\exp(12\gamma^{0,1})=f^*\mathcal O_B(1)$. Its injectivity and the root relation give the stated $M_1$ equality.
\end{definition}
\uses{The standard clockwise affine data and exponential sequence; the root factorization.}

\begin{theorem}[cdp:C5b; T; Integral twists and the normalization-map exception]\label{cdp:C5b}
For $n\in\mathbb Z$ define
\[
 r_3(n)\equiv n\pmod3,\quad 0\le r_3(n)<3,
 \qquad r_4(n)\equiv3n\pmod4,\quad0\le r_4(n)<4,
\]
\begin{gather*}
 d(n)=\frac{n-4r_3(n)-3r_4(n)}{12},\qquad
 \epsilon_n=\begin{cases}0&3\mid n,\\1&3\nmid n.\end{cases}
\\
 u(n)=\max\{d(n)+1,0\},\qquad
 v(n)=\max\{-d(n)-1,0\}.
\end{gather*}
The coarse direct images of $T_X\otimes M_n$ are
\begin{align*}
 f_*(T_X\otimes M_n)&=\mathcal O(d(n))\oplus\mathcal O(d(n-1)),\\
 R^1f_*(T_X\otimes M_n)&=\mathcal O(d(n))\oplus\mathcal O(d(n-1))\oplus\mathcal O(d(n+2)),\\
 R^2f_*(T_X\otimes M_n)&=\mathcal O(d(n+6))\oplus\mathbb C_{p_1}^{\epsilon_n}.\tag{C.5}
\end{align*}
Their global cohomology dimensions are
\begin{align*}
 h^0&=u(n)+u(n-1),\\
 h^1&=v(n)+v(n-1)+u(n)+u(n-1)+u(n+2),\\
 h^2&=v(n)+v(n-1)+v(n+2)+u(n+6)+\epsilon_n,\\
 h^3&=v(n+6).                                             \tag{C.6}
\end{align*}
The natural map $H^2(X,T_X\otimes M_n)\to H^2(\mathrm dP_6,K_{\mathrm dP_6})=\mathbb C$ has rank one exactly for
\[
 n\in\{-6,-3,-2\}\cup\mathbb Z_{\ge0};                         \tag{C.7}
\]
it is zero for all other integers. In particular $M_{-1}$ has tangent cohomology $(0,0,1,0)$ and a zero map to the normalization target.
\end{theorem}
\uses{The root and normalization calculations below; inv:D1.3 for the source-expanded branch-normal and conormal computation; analytic base change \cite[III.4.10, III.4.12]{BS76} and relative duality \cite{RRV71}.}

\begin{proof}
\par\smallskip\noindent\textbf{Step 1.}\ We give the root-chart and cusp calculations used in this proof. The root charts at $p_1,p_2$ are the finite etale torus-family covers used in the construction; finite-group invariants over $\mathbb C$ are exact. At the cusp the local equations are $xy=q$ and $xyz=q$, and the vertical logarithmic fields are generated by $x\partial_x-y\partial_y$ and the remaining relative direction. Their affine transition derivatives depend only on the base. The invariant vertical field gives a trivial subline of the vertical bundle $\mathcal E_{\rm tan}$. The quotient has twelfth power $\mathcal O(-p_0)$: its automorphy is that of the reciprocal elliptic discriminant, whose only divisor is a simple pole at the cusp. As $\Pic(\mathfrak B)=\mathbb Z$, the quotient is $\mathcal O(-1)$; $H^1(\mathfrak B,\mathcal O(1))=0$ splits the extension. Thus $\mathcal E_{\rm tan}=\mathcal O\oplus\mathcal O(-1)$.

\par\smallskip\noindent\textbf{Step 2.}\ The normalization complex of $W$ has global scalar complex
\[
 \mathbb C\xrightarrow{0}\mathbb C^3
 \xrightarrow{(a,b,c)\mapsto(a-b+c,a-b+c)}\mathbb C^2.
\]
The rational normalization and the three rational double curves have no positive-degree structure-sheaf cohomology. The displayed complex therefore gives $(1,2,1)$ for the fibre cohomology and its exterior cup product; the two oriented triangular faces give the nondegenerate degree-one pairing. The cyclic-cover fibres have the same dimensions. Proper flat base change on these charts gives locally free $G_{\mathfrak B}=R^1\bar f_*\mathcal O_X$, with $R^2\bar f_*\mathcal O_X=\det G_{\mathfrak B}$. Relative duality and the determinant of the logarithmic sequence give $\det G_{\mathfrak B}=\mathcal O(-1)$. The integral class $\gamma$ maps to a nowhere-zero section of $G_{\mathfrak B}$: on a marked smooth torus the two identity columns of the period matrix prevent it from belonging to the span of the two holomorphic period rows, and at the cusp it is a nonzero cycle in the displayed complex. The same splitting argument yields
\[
\mathcal E_{\rm tan}=G_{\mathfrak B}=\mathcal O\oplus\mathcal O(-1),\quad R^2\bar f_*\mathcal O_X=\mathcal O(-1),
\quad0\to\bar f^*\mathcal E_{\rm tan}\to\mathcal T\to\bar f^*\mathcal O(-5)\to0. \tag{C.8}
\]
These are the actual root-chart and cusp comparisons used here; no coarse multiple fibre was replaced by its reduction in this argument.

\par\smallskip\noindent\textbf{Step 3.}\ Choose root sections $x$ of weight $3$, vanishing at $p_2$, and $y$ of weight $4$, vanishing at $p_1$. The Kodaira--Spencer morphism $\mathcal O(-5)\to \mathcal E_{\rm tan}\otimes G_{\mathfrak B}$ has, by its weights, the form
\[
 \kappa=\begin{pmatrix}0&a y\\b y&c_0x\end{pmatrix}.
\]
Its first coefficient vanishes because $H^0(\mathcal O(5))=0$. The logarithmic period derivative at the cusp is invertible, giving $ab\ne0$. At $p_1$, the elliptic coordinate $\tau-\rho$ has a simple zero on the root chart, so its period derivative is nonzero, giving $c_0\ne0$. Rescaling the four line frames makes the three nonzero coefficients one. Cup product with this extension class gives
\[
\delta_0=(0,y,y,x)^t,\qquad
\delta_1=\begin{pmatrix}-y&0\\-x&y\end{pmatrix}:
 \mathcal O(-5)\oplus\mathcal O(-6)\to\mathcal O(-1)\oplus\mathcal O(-2).
\]
The first is a subbundle inclusion since $x,y$ have no common zero, and its cokernel is $\mathcal O\oplus\mathcal O(-1)\oplus\mathcal O(2)$. The second is injective with determinant $-y^2$; its cokernel is
$Q=\mathcal O(-1)|_{2p_1}$, a length-two module on the order-three chart. The direct-image sequence of (C.8) yields
\[
 R^0\bar f_*\mathcal T=\mathcal E_{\rm tan},\quad
 R^1\bar f_*\mathcal T=\mathcal O\oplus\mathcal O(-1)\oplus\mathcal O(2),\quad
 R^2\bar f_*\mathcal T=\mathcal O(-6)\oplus Q.
\]
The last splitting follows from $H^1(Q(6))=0$; it is not designated canonical.

\par\smallskip\noindent\textbf{Step 4.}\ The branch-normal sequence is
\[
0\longrightarrow\mathcal T\longrightarrow T_X
\longrightarrow\nu_*K_{\mathrm dP_6}\longrightarrow0.
\]
To verify the quotient, in a chart $W=(z_1\cdots z_r=0)$ the map takes
a vector field to its normal coefficient on each of the $r$ separate
branches. The coefficients lift independently; its kernel is precisely
the vector fields tangent to every branch. Thus the quotient is the
pushforward of the normal line of $\nu$, not $\mathcal O_W(W)$.
The form $d\log x_1\wedge d\log x_2\wedge dq$ becomes a unit times
$dz_1\wedge dz_2\wedge dz_3$ in every unimodular chart and is invariant
under every cusp deck transformation. It trivializes $\nu^*K_X$.
Adjunction for the immersion gives its normal line $K_{\mathrm dP_6}$, proving
the normal-coefficient sequence. Since $\mathrm dP_6$ is rational, its $K_{\mathrm dP_6}$ cohomology is $(0,0,1)$ by
Serre duality. Finite pushforward preserves that cohomology.

\par\smallskip\noindent\textbf{Step 5.}\ It remains to identify the extension at $p_0$ in degree two. We give
a fibre-sensitive calculation. The kernel of
\[
 \Omega_X^1|_W\longrightarrow\nu_*\Omega_{\mathrm dP_6}^1
\]
is $\nu_*\mathcal O_{\mathrm dP_6}\otimes N_W^*$. At $xy=0$ its sections are
$a(y,z)y\,dx+b(x,z)x\,dy$; at $xyz=0$ they are
$a(y,z)yz\,dx+b(x,z)xz\,dy+c(x,y)xy\,dz$.
These identifications glue by the common conormal $dq$.
Since $N_W$ and $K_X|_W$ are trivial and $H^0(P,\Omega_{\mathrm dP_6}^1)=0$, the
Serre-dual space of $H^2(W,T_X|_W)$ is one-dimensional.
Top-degree base change identifies this group with the residue fibre
of $(R^2\bar f_*T_X)_{p_0}$. Its generic rank is one by the displayed logarithmic direct images.
A finite module over the discrete valuation ring $\mathbb C\{q\}$ of rank
one and residue dimension one is free of rank one. Thus the extension
by $\mathbb C_{p_0}$ increases the degree of its line part by $p_0$, without
adding torsion. 
Every $\nu^*M_n$ is trivial because $H^2(X,\mathbb Z)=0$ and
$\operatorname{Pic}^0(\mathrm dP_6)=0$. Thus the same quotient
calculation applies after each integral twist. The elementary extension
of $\mathcal O(-6)$ by $\mathbb C_{p_0}$ is $\mathcal O(6)$,
since $\mathcal O(p_0)=\mathcal O(12)$ on the root curve.
We have
\[
 R^0\bar f_*T_X=\mathcal E_{\rm tan},\quad
 R^1\bar f_*T_X=\mathcal O\oplus\mathcal O(-1)\oplus\mathcal O(2),\quad
 R^2\bar f_*T_X=\mathcal O(6)\oplus Q.                       \tag{C.9}
\]

\par\smallskip\noindent\textbf{Step 6.}\ For a root line, removing its two stabilizer characters gives
$\varpi_*\mathcal O(n)=\mathcal O_B(d(n))$. Averaging the finite stabilizers is exact, so there are no higher direct images of $\rho$. Its sections are equivalently monomials $x^ay^b$ with $3a+4b=n$, and Serre duality gives $h^1\mathcal O(n)=h^0\mathcal O(-n-7)$. At $p_1$, the invariant exponent of $Q(n)$ is $n-1$ modulo three. Its two possible coefficients have orders zero and one, so its coarse length is $\epsilon_n$ and it is killed by the coarse parameter. Apply projection formula to (C.9) and then $\varpi_*$. This proves (C.5), including the support and length of every torsion summand. Leray over $B$ has columns only zero and one, proving (C.6).

\par\smallskip\noindent\textbf{Step 7.}\ The normalization map on $R^2$ is the quotient in
\[
0\longrightarrow\mathcal O(n-6)\longrightarrow\mathcal O(n+6)
 \longrightarrow\mathbb C_{p_0}\longrightarrow0.
\]
Its global-section map is nonzero precisely when a section of $\mathcal O(n+6)$ is nonzero at $p_0$. Both $x,y$ are nonzero at that point, so this occurs precisely when $n+6$ lies in the additive semigroup generated by $3,4$. That semigroup is $\{0,3,4,6,7,8,\ldots\}$, giving (C.7). The $H^1(R^1)$ part maps to zero in the cohomology of a sheaf supported at one point, by the natural Leray filtration. At $n=-1$, (C.5) has only one nonzero cohomological contribution, the reduced $\mathbb C_{p_1}$ summand in degree two, while $H^3(X,\mathcal T\otimes M_{-1})=H^1(\mathcal O(-7))=\mathbb C$. The boundary from the normalization target is an isomorphism. This proves the stated counterexample to a surjection for every twist.
\end{proof}

\subsection{Coprime higher multiplicities}\label{sec:tangent-higher}
The tangent vector occurs in two producing runs, \texttt{PUSH} E.2.1
and \texttt{CDP} E11: \emph{two runs, no second reader}.
The statement below retains the full hypotheses and conclusions of
\texttt{PUSH} E.2.1; the comparison with \texttt{CDP} is limited to
their common higher-multiplicity tangent-cohomology conclusion.


\begin{theorem}[push:E.2.1; T; Tangent cohomology and local miniversality for coprime multiplicities]\label{push:E.2.1}
For every sphere member with multiplicities $(3k,4l)$ and the
admissible logarithmic data of push:E.1 (in particular, $\gcd(3k,4l)=1$), put $z=kl$ and $d=3k+4l$.
On $\mathfrak B_{k,l}=\mathbb P(3k,4l)$ the root lift
$\bar f:X\to\mathfrak B_{k,l}$ satisfies
\[
 \begin{split}
 \mathcal E_{\rm tan}&=\mathcal O\oplus\mathcal O(-z),\\
 R^q\bar f_*\mathcal O_X&=
 \mathcal O,\quad\mathcal O\oplus\mathcal O(-z),\quad\mathcal O(-z),\\
 K_X&=\bar f^*\mathcal O(z-d).
 \end{split}
\]
The tangent cohomology is $(h^0,h^1,h^2,h^3)(X,T_X)=(1,1,1,0)$.
With the groupoid derivative convention
$\kappa(\Phi)=(\partial_c\Phi)\Phi^{-1}$, the class of
$\partial/\partial c$ is $e\otimes[\gamma]$.  The class of
$\partial/\partial u$ is
$-(2\pi i u)^{-1}e\otimes[\gamma]$.  Both are nonzero.
The Kuranishi germ is a smooth reduced line, and the fixed-data
period family is locally miniversal.
\end{theorem}
\uses{The admissible higher-multiplicity data push:E.1; the displayed root-chart and residue computations; \cite[III.4.12]{BS76} and \cite{RRV71}. Two runs, no second reader.}

\begin{proof}
\par\smallskip\noindent\textbf{Step 1.}\ The root object is defined by the actual divisor identities
$f^*(0)=mS_1$ and $f^*(1)=nS_2$, with $m=3k$, $n=4l$.
It gives a holomorphic map to the analytic root stack.  Its base
changes to the two root discs are the actual smooth torus
families; the quotient actions are free.  At the cusp its charts
are $q=z_0z_1z_2$, so the source is smooth and the map is flat,
of relative dimension two, with Gorenstein fibres.  Properness is
checked on the root atlas, where it is properness of the compact
torus families and the original compact cusp filling.

\par\smallskip\noindent\textbf{Step 2.}\ Use weighted coordinates $x,y$ of weights $m,n$, with coarse
coordinate $t=y^m/(y^m-x^n)$.
The two root divisors have bundles $\mathcal O(n)$ and
$\mathcal O(m)$, and the cusp has bundle $\mathcal O(mn)$.
The root-forgetting morphism
\[
 r:\mathfrak B_{k,l}\longrightarrow\mathfrak B_{1,1},
 \qquad (x,y)\longmapsto(x^l,y^k),
\]
uses the group homomorphism $\lambda\mapsto\lambda^{kl}$ on
the weighted quotient.  It is a holomorphic morphism of analytic
stacks over $B$, with local maps $s\mapsto s^k$ and
$s\mapsto s^l$.  In particular $r^*\mathcal O(1)=\mathcal O(z)$.
For completeness, the analytic Picard group of this root curve has
the presentation
$\langle H,R_1,R_2\mid mR_1=H,\ nR_2=H\rangle$:
after subtracting the two stabilizer characters, a line bundle
descends to the coarse projective line.  Its only relations are
the two defining root identities.  The Smith form and
$\gcd(m,n)=1$ identify this group with $\mathbb Z$, with
$H=mn\mathcal O(1)$, $R_1=n\mathcal O(1)$, and
$R_2=m\mathcal O(1)$.  This proves the analytic torsion-free
Picard assertion used below.
Translations have identity vertical derivative and act trivially
on structure-sheaf cohomology.  Thus the vertical bundle computed
from the linear period factors is $r^*(\mathcal O\oplus
\mathcal O(-1))$.  The original determinant calculation can also
be repeated: its quotient line $L$ has
$L^{12}=\mathcal O(-p_0)=\mathcal O(-12z)$, so
$L=\mathcal O(-z)$ because the Picard group is torsion-free.
The fixed vector gives a subline $\mathcal O$; its extension
splits because $H^1(\mathcal O(z))=0$.

\par\smallskip\noindent\textbf{Step 3.}\ The cusp normalization resolution has terms
$\mathcal O_{\mathrm{dP}_6}$, three copies of
$\mathcal O_{\mathbb P^1}$, and two point sheaves.
Its global differential is $(a,b,c)\mapsto
(a-b+c,a-b+c)$, giving cohomology $(1,2,1)$; its alternating cup
pairing in degree one is nondegenerate.  The degree-one cocycle
$\gamma$ extends across the cusp because it vanishes on the toric
period lattice.  Its edge values $(0,-1,-1)$ have nonzero cup
product with $(1,0,-1)$.  On finite root charts its class is the
nonzero constant lattice cocycle $\gamma$, invariant under every
deck transformation.  These are explicit comparisons of the
same class, so it gives a nowhere-zero subline in
$G_{\mathfrak B}=R^1\bar f_*\mathcal O_X$.
The base-change and relative-duality argument applies on these charts, using \cite[III.4.12]{BS76} and \cite[pp.~264--266]{RRV71}.  It gives $R^2\bar f_*\mathcal O_X=L$ and
$\det G_{\mathfrak B}=L$.  The extension of $L$ by the specified
$[\gamma]$-subline splits by the same $H^1(\mathcal O(z))=0$.
This proves all direct-image formulas.

\par\smallskip\noindent\textbf{Step 4.}\ The logarithmic tangent sequence is
\[
 0\longrightarrow\mathcal V=\bar f^*\mathcal E_{\rm tan}
 \longrightarrow\mathcal T=T_X(-\log W)
 \longrightarrow\mathcal H=\bar f^*\mathcal O(d-12z)
 \longrightarrow0.
\]
Its last map is the relative differential on smooth root charts
and sends $\sum a_i z_i\partial_{z_i}$ to
$(\sum a_i)q\partial_q$ at the cusp.  It is surjective there.
Since $K_{\mathfrak B}=\mathcal O(-d)$, determinants give
$K_X=\bar f^*\mathcal O(z-d)$.

\par\smallskip\noindent\textbf{Step 5.}\ We give the weighted cohomology calculation explicitly; simply
replacing $1$ by $kl$ in an acyclicity assertion would be incorrect.
The two quotient charts $x\ne0$, $y\ne0$ are finite quotients
of affine lines, and their overlap is a finite quotient of a
punctured line.  Finite averaging and Stein acyclicity compute
coherent cohomology by this two-chart complex.  Its homogeneous
Laurent monomials give
\[
 \begin{split}
 h^0\mathcal O(a)&=\#\{(i,j)\in\mathbb Z_{\ge0}^2:mi+nj=a\},\\
 h^1\mathcal O(a)&=h^0\mathcal O(-a-d).
 \end{split}
\]
The second equality identifies the remaining monomials
$x^{-i-1}y^{-j-1}$; it does not require a positivity theorem.
The bundles $\mathcal O(-z),\mathcal O(-2z)$ are acyclic.
Indeed a possible degree-one monomial would give
$m(i+1)+n(j+1)=az$, with $a=1$ or $2$.
Coprimality forces $i+1=lA$, $j+1=kB$, hence
$3A+4B=a$ with $A,B\ge1$, which is impossible.
Positive degree has no degree-one monomials, so
$H^1\mathcal O(z)=0$ as used above.

\par\smallskip\noindent\textbf{Step 6.}\ Put $b=d-12z$.  Both $b$ and $b-z$ are negative.
Degree-one monomials for $\mathcal O(b)$ would satisfy
$3A+4B=12$, with $A,B\ge1$, and none exist.
For $\mathcal O(b-z)$ the equation is $3A+4B=13$;
its only positive solution is $(A,B)=(3,1)$, producing
$x^{3l-2}y^{k-2}$ in the dual description.  It is allowed
exactly when $k>1$; write $e_k=1$ in that case and $e_k=0$
otherwise.  Projection formula and the two-column Leray sequence
now give
\[
 h^\bullet(\mathcal V)=(1,1,0,0),\qquad
 h^\bullet(\mathcal H)=(0,0,e_k,e_k),\qquad
 h^\bullet(\mathcal T)=(1,1,e_k,e_k).
\]
The producing run reports that \texttt{e\_push\_exact.py} independently checks these
counts for all $1965$ coprime pairs with $k,l\le80$; the monomial
argument proves them for all pairs.

\par\smallskip\noindent\textbf{Step 7.}\ On the normalization $\nu:P_6=\mathrm{dP}_6\to W\subset X$,
the branchwise normal line is $K_{P_6}$: the invariant cusp form
$d\log x_1\wedge d\log x_2\wedge dq$ trivializes $K_X$
on the entire cusp neighbourhood.  Local normal coefficients
therefore give the exact sequence
\[
 0\longrightarrow\mathcal T\longrightarrow T_X
 \longrightarrow\nu_*K_{P_6}\longrightarrow0,
 \qquad h^\bullet(K_{P_6})=(0,0,1).
\]
We must determine its connecting map when $k>1$.
The dual logarithmic sequence, tensored with $K_X$, has horizontal
subline $\bar f^*\mathcal O(13z-2d)$ and vertical quotient
$\bar f^*(\mathcal O(z-d)\oplus\mathcal O(2z-d))$.
The latter has no global sections, by the same equations
$3A+4B=1,2$ above.  Consequently
\[
 H^0(\Omega_X^1(\log W)\otimes K_X)
 =H^0(\mathfrak B,\mathcal O(13z-2d)).
\]
For $k>1$ this is generated by the monomial
$x^{3l-2}y^{k-2}$ multiplying the horizontal logarithmic
differential.  Both $x$ and $y$ are nonzero over $p_0$, so its
residue on $P_6$ is nonzero.  The ordinary logarithmic residue
sequence, tensored with $K_X$, identifies this residue map with
the Serre dual of
$H^2(P_6,K_{P_6})\to H^3(X,\mathcal T)$.
One can check this identification directly in each chart
$q=z_0z_1z_2$: the residues of $dq/q$ are one on each branch,
and contraction with the invariant three-form is the normal
coefficient pairing.  Thus that connecting map has rank $e_k$.
The long exact sequence proves $h^2(T_X)=1$ and $h^3(T_X)=0$;
degrees zero and one are unchanged.

\par\smallskip\noindent\textbf{Step 8.}\ Finally $\partial_c\Pi=e\otimes\gamma$ gives the constant
groupoid cocycle $\lambda\mapsto\gamma(\lambda)e$ on every
regular torus.  On the enlarged finite roots the derivative of the
affine generator is $(r_j/m_j)e$, and the derivative of its
punctured connector is $(\log s_j)/(2\pi i)\,r_je$.
The identity $m_j(r_j/m_j)=\gamma(w_j)$ verifies the power
relation, while $\gamma A_j=\gamma$ verifies conjugation with
every lattice arrow.  Changing $\log s_j$ by $2\pi i N$
adds the derivative of the integral arrow $Nw_j$; the connector
formula is used only on the punctured overlap.  Its central
root-chart representative is the constant $(r_j/m_j)e$,
which is holomorphic.  At the cusp $\gamma$ kills the toric periods,
so the constant deck cocycle is exactly its extending cocycle
above.  Its image under the two isomorphisms
$H^1(\mathcal V)\simeq H^1(\mathcal T)\simeq H^1(T_X)$
is $e\otimes[\gamma]\ne0$.  The chain rule gives the asserted
$u$-class.  The Kuranishi local algebra is
$\mathbb C\{t\}/I$ since $h^1=1$.  The existing family induces
$t\mapsto g(c)$ with $g'(c_0)\ne0$.  The analytic inverse
function theorem makes this substitution injective; every
element of $I$ maps to zero, so $I=0$.  Its classifying map is
then an isomorphism of germs, proving local miniversality.
\end{proof}

\subsection{The algebraic-dimension corollary}\label{sec:algebraic-dimension-nearby}

\begin{corollary}[fam:A.5.2-corollary; T; Algebraic dimension near the standard family]\label{fam:A.5.2-corollary}
Every sufficiently small deformation of $X_c$ has algebraic dimension $1$.
\end{corollary}
\uses{fam:A.5.2; \cite[Theorem 9.1(i)]{AC26}.}

\begin{proof}
By fam:A.5.2, every sufficiently small deformation is induced by the compact period family. Each member has algebraic dimension one by \cite[Theorem 9.1(i)]{AC26}. This proves the stated local corollary.
\end{proof}
